%% Research in Engineering Design -- Springer Nature template (sn-jnl), author-year references.
%DIF LATEXDIFF DIFFERENCE FILE
%DIF DEL old.tex    Tue Oct  6 17:27:32 2026
%DIF ADD main.tex   Tue Oct  6 18:09:02 2026
%% Line numbers for review (option lineno); remove the option for the final version.
\documentclass[pdflatex,sn-basic,lineno]{sn-jnl}

\usepackage{graphicx}
\usepackage{multirow}
\usepackage{amsmath,amssymb,amsfonts}
\usepackage{amsthm}
\usepackage{booktabs}
\usepackage{array}
\usepackage{tikz}
\usetikzlibrary{arrows.meta,positioning,calc}
\usepackage{siunitx}
\usepackage{manyfoot}   % the class's begin-document footnote hook expects this package to be loaded
\usepackage{placeins}   % \FloatBarrier keeps floats from drifting past sections
\graphicspath{{figures/}}

\raggedbottom

%% Float placement: allow figures and tables to share pages with text and with each other.
\renewcommand{\topfraction}{0.9}
\renewcommand{\bottomfraction}{0.7}
\renewcommand{\textfraction}{0.07}
\renewcommand{\floatpagefraction}{0.8}
\setcounter{topnumber}{3}
\setcounter{totalnumber}{4}
\setlength{\textfloatsep}{10pt plus 3pt minus 2pt}
\setlength{\floatsep}{10pt plus 3pt minus 2pt}
\setlength{\intextsep}{10pt plus 3pt minus 2pt}
\setlength{\belowcaptionskip}{3pt}
%% Float pages fill from the top instead of centring their floats.
\makeatletter
\setlength{\@fptop}{0pt}
\setlength{\@fpsep}{12pt plus 2fil}
\setlength{\@fpbot}{0pt plus 1fil}
\makeatother
%% Table body size: the class's own 8 bp table font (its \footnotesize is 7 pt and its \scriptsize 9 pt).
\newcommand{\tabsize}{\fontsize{8}{9.5}\selectfont}
%DIF 39-40c39-46
%DIF < %% Reference list: the class separates entries by 1em; half of that keeps them distinct on fewer pages.
%DIF < \setlength{\bibsep}{0.5em}
%DIF -------
%% Reference list: the class sets it at normal size with entries 1em apart; the 8 bp table font with %DIF > 
%% hanging indents keeps the entries distinct on fewer pages. %DIF > 
\setlength{\bibsep}{0pt} %DIF > 
\makeatletter %DIF > 
\def\bibfont{\reset@font\fontfamily{\rmdefault}\fontsize{8}{9.5}\selectfont} %DIF > 
\makeatother %DIF > 
%% DOIs in the reference list as doi:..., linked to doi.org. %DIF > 
\newcommand{\doi}[1]{\href{https://doi.org/#1}{doi:\nolinkurl{#1}}} %DIF > 
%DIF -------
%% Ragged-right paragraph column for text tables.
\newcolumntype{L}[1]{>{\raggedright\arraybackslash\hspace{0pt}}p{#1}}

%% Set \anontrue to produce an anonymised copy (no author block, no identifying declarations).
\newcommand{\pending}[1]{\textit{[#1 pending]}}
\newif\ifanon
\anonfalse
%DIF PREAMBLE EXTENSION ADDED BY LATEXDIFF
%DIF UNDERLINE PREAMBLE %DIF PREAMBLE
\RequirePackage[normalem]{ulem} %DIF PREAMBLE
\RequirePackage{color}\definecolor{RED}{rgb}{1,0,0}\definecolor{BLUE}{rgb}{0,0,1} %DIF PREAMBLE
\providecommand{\DIFadd}[1]{{\protect\color{blue}\uwave{#1}}} %DIF PREAMBLE
\providecommand{\DIFdel}[1]{{\protect\color{red}\sout{#1}}}                      %DIF PREAMBLE
%DIF SAFE PREAMBLE %DIF PREAMBLE
\providecommand{\DIFaddbegin}{} %DIF PREAMBLE
\providecommand{\DIFaddend}{} %DIF PREAMBLE
\providecommand{\DIFdelbegin}{} %DIF PREAMBLE
\providecommand{\DIFdelend}{} %DIF PREAMBLE
\providecommand{\DIFmodbegin}{} %DIF PREAMBLE
\providecommand{\DIFmodend}{} %DIF PREAMBLE
%DIF FLOATSAFE PREAMBLE %DIF PREAMBLE
\providecommand{\DIFaddFL}[1]{\DIFadd{#1}} %DIF PREAMBLE
\providecommand{\DIFdelFL}[1]{\DIFdel{#1}} %DIF PREAMBLE
\providecommand{\DIFaddbeginFL}{} %DIF PREAMBLE
\providecommand{\DIFaddendFL}{} %DIF PREAMBLE
\providecommand{\DIFdelbeginFL}{} %DIF PREAMBLE
\providecommand{\DIFdelendFL}{} %DIF PREAMBLE
\newcommand{\DIFscaledelfig}{0.5}
%DIF HIGHLIGHTGRAPHICS PREAMBLE %DIF PREAMBLE
\RequirePackage{settobox} %DIF PREAMBLE
\RequirePackage{letltxmacro} %DIF PREAMBLE
\newsavebox{\DIFdelgraphicsbox} %DIF PREAMBLE
\newlength{\DIFdelgraphicswidth} %DIF PREAMBLE
\newlength{\DIFdelgraphicsheight} %DIF PREAMBLE
% store original definition of \includegraphics %DIF PREAMBLE
\LetLtxMacro{\DIFOincludegraphics}{\includegraphics} %DIF PREAMBLE
\newcommand{\DIFaddincludegraphics}[2][]{{\color{blue}\fbox{\DIFOincludegraphics[#1]{#2}}}} %DIF PREAMBLE
\newcommand{\DIFdelincludegraphics}[2][]{% %DIF PREAMBLE
\sbox{\DIFdelgraphicsbox}{\DIFOincludegraphics[#1]{#2}}% %DIF PREAMBLE
\settoboxwidth{\DIFdelgraphicswidth}{\DIFdelgraphicsbox} %DIF PREAMBLE
\settoboxtotalheight{\DIFdelgraphicsheight}{\DIFdelgraphicsbox} %DIF PREAMBLE
\scalebox{\DIFscaledelfig}{% %DIF PREAMBLE
\parbox[b]{\DIFdelgraphicswidth}{\usebox{\DIFdelgraphicsbox}\\[-\baselineskip] \rule{\DIFdelgraphicswidth}{0em}}\llap{\resizebox{\DIFdelgraphicswidth}{\DIFdelgraphicsheight}{% %DIF PREAMBLE
\setlength{\unitlength}{\DIFdelgraphicswidth}% %DIF PREAMBLE
\begin{picture}(1,1)% %DIF PREAMBLE
\thicklines\linethickness{2pt} %DIF PREAMBLE
{\color[rgb]{1,0,0}\put(0,0){\framebox(1,1){}}}% %DIF PREAMBLE
{\color[rgb]{1,0,0}\put(0,0){\line( 1,1){1}}}% %DIF PREAMBLE
{\color[rgb]{1,0,0}\put(0,1){\line(1,-1){1}}}% %DIF PREAMBLE
\end{picture}% %DIF PREAMBLE
}\hspace*{3pt}}} %DIF PREAMBLE
} %DIF PREAMBLE
\LetLtxMacro{\DIFOaddbegin}{\DIFaddbegin} %DIF PREAMBLE
\LetLtxMacro{\DIFOaddend}{\DIFaddend} %DIF PREAMBLE
\LetLtxMacro{\DIFOdelbegin}{\DIFdelbegin} %DIF PREAMBLE
\LetLtxMacro{\DIFOdelend}{\DIFdelend} %DIF PREAMBLE
\DeclareRobustCommand{\DIFaddbegin}{\DIFOaddbegin \let\includegraphics\DIFaddincludegraphics} %DIF PREAMBLE
\DeclareRobustCommand{\DIFaddend}{\DIFOaddend \let\includegraphics\DIFOincludegraphics} %DIF PREAMBLE
\DeclareRobustCommand{\DIFdelbegin}{\DIFOdelbegin \let\includegraphics\DIFdelincludegraphics} %DIF PREAMBLE
\DeclareRobustCommand{\DIFdelend}{\DIFOaddend \let\includegraphics\DIFOincludegraphics} %DIF PREAMBLE
\LetLtxMacro{\DIFOaddbeginFL}{\DIFaddbeginFL} %DIF PREAMBLE
\LetLtxMacro{\DIFOaddendFL}{\DIFaddendFL} %DIF PREAMBLE
\LetLtxMacro{\DIFOdelbeginFL}{\DIFdelbeginFL} %DIF PREAMBLE
\LetLtxMacro{\DIFOdelendFL}{\DIFdelendFL} %DIF PREAMBLE
\DeclareRobustCommand{\DIFaddbeginFL}{\DIFOaddbeginFL \let\includegraphics\DIFaddincludegraphics} %DIF PREAMBLE
\DeclareRobustCommand{\DIFaddendFL}{\DIFOaddendFL \let\includegraphics\DIFOincludegraphics} %DIF PREAMBLE
\DeclareRobustCommand{\DIFdelbeginFL}{\DIFOdelbeginFL \let\includegraphics\DIFdelincludegraphics} %DIF PREAMBLE
\DeclareRobustCommand{\DIFdelendFL}{\DIFOaddendFL \let\includegraphics\DIFOincludegraphics} %DIF PREAMBLE
%DIF COLORLISTINGS PREAMBLE %DIF PREAMBLE
\RequirePackage{listings} %DIF PREAMBLE
\RequirePackage{color} %DIF PREAMBLE
\lstdefinelanguage{DIFcode}{ %DIF PREAMBLE
%DIF DIFCODE_UNDERLINE %DIF PREAMBLE
  moredelim=[il][\color{red}\sout]{\%DIF\ <\ }, %DIF PREAMBLE
  moredelim=[il][\color{blue}\uwave]{\%DIF\ >\ } %DIF PREAMBLE
} %DIF PREAMBLE
\lstdefinestyle{DIFverbatimstyle}{ %DIF PREAMBLE
	language=DIFcode, %DIF PREAMBLE
	basicstyle=\ttfamily, %DIF PREAMBLE
	columns=fullflexible, %DIF PREAMBLE
	keepspaces=true %DIF PREAMBLE
} %DIF PREAMBLE
\lstnewenvironment{DIFverbatim}{\lstset{style=DIFverbatimstyle}}{} %DIF PREAMBLE
\lstnewenvironment{DIFverbatim*}{\lstset{style=DIFverbatimstyle,showspaces=true}}{} %DIF PREAMBLE
%DIF END PREAMBLE EXTENSION ADDED BY LATEXDIFF

\begin{document}

\title[Evaluating design-stage manufacturability decisions]{Evaluating design-stage manufacturability decisions against the resolution of production}

\ifanon
\author*[1]{\fnm{Anonymised} \sur{for review}}
\affil*[1]{\orgname{Affiliation withheld for review}}
\else
\author*[1]{\fnm{H\"useyin Tayyer} \sur{Canseven}}\email{huseyin.canseven@lut.fi}

\affil*[1]{\orgdiv{Department of Electrical Engineering, LUT School of Energy Systems}, \orgname{LUT University}, \orgaddress{\street{Yliopistonkatu 34}, \city{Lappeenranta}, \postcode{53850}, \country{Finland}}}
\fi

%% BEGIN ABSTRACT
\DIFdelbegin %DIFDELCMD < \abstract{Design-stage manufacturability decisions increasingly rest on simulation and machine learning, yet whether a prediction is fit to decide is rarely evaluated against production. A framework is proposed whose unit, the production floor, is the difference that drift and scatter produce between two batches of one design within a run. A rule returning meets, fails or trial is characterised by a decisive distance in floors, beyond which at least 95\,\% of its verdicts are correct, and a safe distance, beyond which at most 5\,\% are wrong, for each relation of a new design to the produced evidence: a variant, a new process setting or a new family. A procedure turns them into a guard band for single verdicts. Thirteen rules were evaluated at a 95\,\% conformance level on open data of 18 deep-drawn design--process alternatives of 500 parts each, mostly process settings on one tool, with matched simulations; one floor of draw-in is 0.05\,mm. With a produced sibling the nearest produced setting decided within 3.4 floors; for a new process setting no rule decided closer than 8.5 floors; in two transfers to a new family none was safe closer than 16. Within a family the choice of rule mattered more than the relation, and requirements at a performance index of 1.33 lay inside every decisive distance. With six to nine produced settings per fit, the simulation, as used, helped only across families, and learned models added abstention rather than accuracy. The framework gives evaluations of AI-based manufacturability support a production-referenced standard.}
%DIFDELCMD < %%%
\DIFdelend \DIFaddbegin \abstract{Design-stage manufacturability decisions rest increasingly on simulation and machine learning, whose fitness to decide is rarely evaluated against production. A framework is proposed whose unit, the production floor, is the difference that drift and scatter of production and measurement produce between two batches of one design within a run. For each relation of a new design to produced evidence, a rule returning meets, fails or trial has a decisive distance in floors, beyond which at least 95\,\% of its verdicts are correct, and a safe distance, beyond which at most 5\,\% are wrong; a procedure adds guard bands for single verdicts. Fourteen rules were evaluated at 95\,\% conformance on open data of 18 deep-drawn alternatives of 500 parts, mostly process settings, with matched simulations; a draw-in floor is 0.05\,mm. With a produced sibling, the nearest produced setting was right for requirements over 3.4 floors from the truth (0.85 for near-replicates, 4.9 for resolvable siblings). For a new setting the best rules needed 4.7 floors interpolating, 4.9 extrapolating to 500\,kN and 8.9 to 100\,kN, with a stroke-speed change; for a new family the best was safe at 16--17 floors in either family's floor. Release-level requirements lay inside every decisive distance except a near-replicate's. With six to nine produced alternatives at two or three forces, the simulation as used, offset or rescaled, helped only across families, and learned 90\,\% intervals covered 85--96\,\% with a sibling, under 50\,\% extrapolating or across families. The framework gives AI-based manufacturability support a production-referenced evaluation standard.}
\DIFaddend %% END ABSTRACT

\keywords{design for manufacturing, machine learning, design decision making, verification and validation, operating characteristic, deep drawing}

\maketitle

%% ----------------------------------------------------------------------
\section{Introduction}\label{sec:intro}

Whether a design alternative can be made to its requirements is decided long before it is made. In sheet-metal forming \DIFdelbegin \DIFdel{, as in most processes, }\DIFdelend the decision rests on a finite-element simulation \citep{banabic2010sheet}: the team compares the predicted characteristics with the requirements and releases the alternative, changes it or sends it to a physical trial. Simulation has moved such decisions forward \citep{thomke1998simulation}, and machine learning makes the loop faster: surrogates replace simulations, models learn manufacturability from earlier designs, and checks are automated \citep{yan2026machine, ghadai2018learning, zhang2018featurenet}. Neither has changed the question the team has to answer: \emph{can this prediction be relied on to decide?} \DIFdelbegin \DIFdel{A simulation that is right on average can be wrong for the alternative at hand, and a model learned on earlier products may be asked about a design unlike any of them. }\DIFdelend Whitney argued in this journal that mechanical design cannot be verified by simulation the way integrated circuits are \citep{whitney1996why}\DIFdelbegin \DIFdel{; }\DIFdelend \DIFaddbegin \DIFadd{, and }\DIFaddend production variation, with which development organisations still contend \citep{thornton2000more}, adds a further reason\DIFdelbegin \DIFdel{. Whether }\DIFdelend \DIFaddbegin \DIFadd{: whether }\DIFaddend a simulation-based decision is credible is \DIFdelbegin \DIFdel{in the end }\DIFdelend a property of production, and production evidence is \DIFdelbegin \DIFdel{what is }\DIFdelend scarce at the design stage.

This paper treats the design-stage manufacturability decision itself as the object of study. How can the fitness of a design-stage prediction \DIFdelbegin \DIFdel{for such a decision }\DIFdelend \DIFaddbegin \DIFadd{to decide }\DIFaddend be expressed on a common production-referenced scale, for any predictor and characteristic (RQ1)? How does it depend on the amount of production evidence and on how the new design relates to what has been produced (RQ2)? And what do simulation and machine learning contribute once production evidence exists (RQ3)? Decision-based design \DIFdelbegin \DIFdel{\mbox{%DIFAUXCMD
\citep{hazelrigg1998framework, lewis2006decision}}\hskip0pt%DIFAUXCMD
}\DIFdelend \DIFaddbegin \DIFadd{\mbox{%DIFAUXCMD
\citep{hazelrigg1998framework}}\hskip0pt%DIFAUXCMD
}\DIFaddend , design margins \DIFdelbegin \DIFdel{\mbox{%DIFAUXCMD
\citep{eckert2019design, brahma2020margin} }\hskip0pt%DIFAUXCMD
}\DIFdelend \DIFaddbegin \DIFadd{\mbox{%DIFAUXCMD
\citep{eckert2019design} }\hskip0pt%DIFAUXCMD
}\DIFaddend and the planning of verification and testing \citep{thomke2001sequential, salado2018mathematical, tahera2019testing} all need to know how far a prediction can be relied on\DIFdelbegin \DIFdel{before a test is cheaper than a decision}\DIFdelend .

The framework proposed here has the \emph{production floor} as its unit: the difference that drift and scatter produce between two batches of \DIFdelbegin \DIFdel{consecutive parts of }\DIFdelend one design within a run, a floor in the sense of a noise floor\DIFdelbegin \DIFdel{rather than of the shop floor}\DIFdelend . A rule that turns the design-stage evidence into a verdict (the alternative \emph{meets} the requirement, \emph{fails} it, or goes to a physical \emph{trial}) is characterised\DIFdelbegin \DIFdel{by two operating characteristics in floors, its }\emph{\DIFdel{decisive distance}} %DIFAUXCMD
\DIFdel{and its }\emph{\DIFdel{safe distance}}%DIFAUXCMD
\DIFdelend , for each \emph{evidence relation} between a new design and the produced alternatives, \DIFaddbegin \DIFadd{by two points on its operating characteristics in floors, its }\emph{\DIFadd{decisive}} \DIFadd{and }\emph{\DIFadd{safe distances}}\DIFadd{, }\DIFaddend and a procedure turns them into a guard band for single verdicts. The decisions in scope are those of a team that has produced variants of a design family, in process planning and tryout and when production knowledge is carried to a new family\DIFdelbegin \DIFdel{; early product design is covered only by the new-family relation}\DIFdelend .

The evaluation uses open data: 9,000 deep-drawn and cut cups of dual-phase steel, produced as 18 design--process alternatives of 500 consecutive parts each \citep{baum2026rddac}, with finite-element simulations of the same process \citep{baum2025ddacs}. Within each of two cup geometries the alternatives differ in blank-holder force and lubrication on one tool, so two of the three relations concern process settings, and the geometries give two transfers between families. Each alternative is judged as a new design \DIFdelbegin \DIFdel{, and the verdicts of thirteen }\DIFdelend \DIFaddbegin \DIFadd{by fourteen }\DIFaddend rules at a 95\,\% conformance level\DIFdelbegin \DIFdel{are compared with }\DIFdelend \DIFaddbegin \DIFadd{, against }\DIFaddend what its 500 parts did. In the terms of the design research methodology \citep{blessing2009drm}, this is a prescriptive study with an initial support evaluation: it shows whether the framework discriminates between rules and evidence situations, not whether design teams decide better with it, and its findings on one dataset are hypotheses for confirmation. \DIFdelbegin %DIFDELCMD < 

%DIFDELCMD < %%%
\DIFdel{The }\DIFdelend \DIFaddbegin \DIFadd{After the positioning (Section~\ref{sec:background}), the }\DIFaddend contributions are the framework (Section~\ref{sec:method})\DIFdelbegin \DIFdel{; }\DIFdelend \DIFaddbegin \DIFadd{, }\DIFaddend an evaluation design that separates the evidence relations \DIFdelbegin \DIFdel{by grouped cross-validation \mbox{%DIFAUXCMD
\citep{roberts2017cross} }\hskip0pt%DIFAUXCMD
}\DIFdelend and isolates what the simulation adds (Section~\ref{sec:demo})\DIFdelbegin \DIFdel{; }\DIFdelend \DIFaddbegin \DIFadd{, }\DIFaddend the empirical answers for deep drawing (Section~\ref{sec:results})\DIFdelbegin \DIFdel{; }\DIFdelend \DIFaddbegin \DIFadd{, }\DIFaddend and guidance and a reporting standard for evaluating simulation- and AI-based manufacturability support (Section~\ref{sec:discussion}).
\DIFdelbegin \DIFdel{Section~\ref{sec:background} positions the work, Section~\ref{sec:conclusion} concludes, and Appendix~\ref{app:appendix} gives the notation, the robustness results and computation details.
}\DIFdelend 

%% ----------------------------------------------------------------------
\section{Background and related work}\label{sec:background}

\subsection{Manufacturability evaluation and machine learning}

Design for manufacturing (DFM) codifies what a process can make into rules, guidelines and cost models \citep{boothroyd2010product}, and quantitative manufacturability models have \DIFdelbegin \DIFdel{been used }\DIFdelend \DIFaddbegin \DIFadd{driven design decisions }\DIFaddend in this journal \DIFdelbegin \DIFdel{to drive design decisions }\DIFdelend \citep{latrobe2003design}. Machine learning learns manufacturability from CAD models \DIFdelbegin \DIFdel{\mbox{%DIFAUXCMD
\citep{ghadai2018learning, zhang2018featurenet, williams2019design}}\hskip0pt%DIFAUXCMD
}\DIFdelend \DIFaddbegin \DIFadd{\mbox{%DIFAUXCMD
\citep{ghadai2018learning, zhang2018featurenet}}\hskip0pt%DIFAUXCMD
}\DIFaddend , from process data and simulated design spaces \citep{yan2026machine}, and generative models propose designs directly \citep{regenwetter2022deep}\DIFdelbegin \DIFdel{; transfer between processes and product families remains difficult \mbox{%DIFAUXCMD
\citep{safdar2024transferability}}\hskip0pt%DIFAUXCMD
}\DIFdelend . Learned manufacturability models are \DIFdelbegin \DIFdel{typically }\DIFdelend evaluated by held-out accuracy on labelled designs\DIFdelbegin \DIFdel{\mbox{%DIFAUXCMD
\citep{ghadai2018learning, zhang2018featurenet, williams2019design}}\hskip0pt%DIFAUXCMD
}\DIFdelend \DIFaddbegin \DIFadd{, and language and vision-language models by case studies \mbox{%DIFAUXCMD
\citep{makatura2024how} }\hskip0pt%DIFAUXCMD
and benchmark questions \mbox{%DIFAUXCMD
\citep{picard2025concept}}\hskip0pt%DIFAUXCMD
}\DIFaddend , not by the decisions they support against the production of the alternatives they are asked about.

\subsection{Design decisions, design types and production evidence}

Decision-based design casts design as a choice between alternatives with uncertain outcomes \DIFdelbegin \DIFdel{\mbox{%DIFAUXCMD
\citep{hazelrigg1998framework, lewis2006decision}}\hskip0pt%DIFAUXCMD
}\DIFdelend \DIFaddbegin \DIFadd{\mbox{%DIFAUXCMD
\citep{hazelrigg1998framework}}\hskip0pt%DIFAUXCMD
}\DIFaddend . How much a team knows about a new alternative depends on the kind of design, variant, adaptive or original \citep{pahl2007engineering}, organised across products by \DIFdelbegin \DIFdel{product-family and platform design \mbox{%DIFAUXCMD
\citep{simpson2006product, jiao2007product}}\hskip0pt%DIFAUXCMD
. The }\DIFdelend \DIFaddbegin \DIFadd{platform design \mbox{%DIFAUXCMD
\citep{simpson2006product}}\hskip0pt%DIFAUXCMD
; the }\DIFaddend evidence relations of this paper are their counterparts for production evidence\DIFdelbegin \DIFdel{: a new variant of a produced setting, a new process setting and a new family}\DIFdelend \DIFaddbegin \DIFadd{, the design type concerning the novelty of the solution and the evidence relation the novelty of its realisation relative to what has been produced}\DIFaddend . Production knowledge enters design through process-capability data \citep{thornton2004variation, tata1999process}\DIFdelbegin \DIFdel{and tolerance analysis \mbox{%DIFAUXCMD
\citep{chase1991survey}}\hskip0pt%DIFAUXCMD
}\DIFdelend . Unresolved uncertainty is absorbed by margins \DIFdelbegin \DIFdel{\mbox{%DIFAUXCMD
\citep{eckert2019design, brahma2020margin}}\hskip0pt%DIFAUXCMD
, and the part of a margin }\DIFdelend \DIFaddbegin \DIFadd{\mbox{%DIFAUXCMD
\citep{eckert2019design}}\hskip0pt%DIFAUXCMD
, whose part }\DIFaddend that covers uncertainty is distinct from the part that absorbs change \citep{eckert2017safety}. Tests and prototypes, the alternative to a margin, are modelled as a trade between the costs of tests and of errors \citep{thomke1998simulation, thomke2001sequential, salado2018mathematical} and studied as design activities \DIFdelbegin \DIFdel{\mbox{%DIFAUXCMD
\citep{tahera2019testing, camburn2017design}}\hskip0pt%DIFAUXCMD
}\DIFdelend \DIFaddbegin \DIFadd{\mbox{%DIFAUXCMD
\citep{tahera2019testing}}\hskip0pt%DIFAUXCMD
}\DIFaddend . When upstream information is fit to act on is the question of overlapped development \citep{krishnan1997model, terwiesch2002exchanging}, and set-based concurrent engineering codifies production knowledge as feasibility limits \DIFdelbegin \DIFdel{for new designs }\DIFdelend \citep{sobek1999toyota}. All of these assume a measure of how far a design-stage prediction can be relied on; the framework supplies one in a unit set by production.

\subsection{Simulation credibility, calibration and validation}

Verification and validation ask whether a model is accurate enough for its intended use\DIFdelbegin \DIFdel{, and }\DIFdelend \DIFaddbegin \DIFadd{; }\DIFaddend the distinction between its validation and application domains \citep{oberkampf2010verification} is the question that the evidence relation asks of a design\DIFdelbegin \DIFdel{; }\DIFdelend \DIFaddbegin \DIFadd{, and }\DIFaddend risk-informed credibility assessment scales the required evidence with the consequence of a wrong decision \citep{asme2018assessing}. Bayesian calibration adds a discrepancy model and calibration parameters to a simulator \citep{kennedy2001bayesian}, multi-fidelity models scale a cheap model to an expensive one \citep{kennedy2000predicting}, and \DIFdelbegin \DIFdel{omitting the discrepancy biases the parameters \mbox{%DIFAUXCMD
\citep{brynjarsdottir2014learning}}\hskip0pt%DIFAUXCMD
; where }\DIFdelend a model may be relied on \DIFdelbegin \DIFdel{at all is the question of }\DIFdelend \DIFaddbegin \DIFadd{only within }\DIFaddend its valid input domain \citep{malak2010using}. Conformal prediction gives distribution-free intervals \DIFdelbegin \DIFdel{with finite-sample coverage }\DIFdelend under exchangeability \citep{lei2018distribution, barber2021predictive}, also under covariate shift \citep{tibshirani2019conformal} and for \DIFdelbegin \DIFdel{the control of }\DIFdelend decision risks \citep{angelopoulos2024conformal}. In sheet-metal forming, friction \citep{hol2012advanced} and the unloading behaviour of the material \citep{yoshida2002model} are \DIFdelbegin \DIFdel{among the }\DIFdelend main sources of simulation error\DIFdelbegin \DIFdel{; }\DIFdelend \DIFaddbegin \DIFadd{, }\DIFaddend simulation datasets of deep drawing support surrogate and transfer learning \citep{baum2025ddacs, heinzelmann2025comprehensive}, and the deviation between simulation and parts has been decomposed statistically \citep{baum2026statistical}. How close to a requirement such intervals and accuracy statements can be relied on to decide, and how much production evidence that takes, remains open.

\subsection{Decisions with abstention and their operating characteristics}

In conformity assessment a measured value proves conformance, proves non-conformance or leaves the decision open, with guard bands set by the measurement uncertainty and the risks of false acceptance and rejection as the measures of a decision rule \DIFdelbegin \DIFdel{\mbox{%DIFAUXCMD
\citep{iso2017decision, jcgm2012evaluation}}\hskip0pt%DIFAUXCMD
}\DIFdelend \DIFaddbegin \DIFadd{\mbox{%DIFAUXCMD
\citep{iso2017decision}}\hskip0pt%DIFAUXCMD
}\DIFaddend . In selection procedures and acceptance sampling the operating characteristic gives the probability of a decision against the true quality, with an indifference zone inside which no decision is required \citep{bechhofer1954single}. A classifier with a reject option trades errors against abstentions \citep{chow1970optimum, geifman2017selective}, learning to defer passes difficult cases to an expert \citep{madras2018predict}, and appropriate reliance on automation requires that its reliability be visible to its users \citep{lee2004trust}, also in design teams working with AI \citep{zhang2021cautionary}. The decisive and safe distances are operating characteristics of this kind, and the \DIFdelbegin \DIFdel{guard band of the procedure }\DIFdelend \DIFaddbegin \DIFadd{procedure's guard band }\DIFaddend is the design-stage counterpart of a conformity-assessment \DIFdelbegin \DIFdel{guard band}\DIFdelend \DIFaddbegin \DIFadd{one}\DIFaddend .

\subsection{Production variation and measurement}

Statistical quality control compares short- and long-term variation with tolerances through capability and performance indices \citep{montgomery2019introduction}, with release criteria such as a performance index of at least 1.33 \citep{aiag2006production} and time-dependent models for drifting processes \citep{iso2026process}. Precision experiments define the repeatability limit, the difference that two results under repeatability conditions exceed with a probability of 5\,\% \citep{iso2025repeatability}, and \DIFaddbegin \DIFadd{intermediate precision under conditions that differ in time, operator or equipment \mbox{%DIFAUXCMD
\citep{iso2023intermediate}}\hskip0pt%DIFAUXCMD
; }\DIFaddend measurement systems analysis separates repeatability from stability \citep{aiag2010measurement}. The production floor is a \DIFdelbegin \DIFdel{repeatability-type limit of batch centres under within-run conditions}\DIFdelend \DIFaddbegin \DIFadd{precision limit of this kind for batch centres within a run (Section~\ref{sec:floor})}\DIFaddend ; what is new is its use \DIFaddbegin \DIFadd{not to judge whether two results agree but }\DIFaddend as the unit \DIFdelbegin \DIFdel{of }\DIFdelend \DIFaddbegin \DIFadd{in which }\DIFaddend a design decision rule \DIFaddbegin \DIFadd{is scored and its guard band set}\DIFaddend . Table~\ref{tab:positioning} places the framework among \DIFdelbegin \DIFdel{these approaches}\DIFdelend \DIFaddbegin \DIFadd{them}\DIFaddend : it does not replace calibration or conformal prediction but scores them \DIFdelbegin \DIFdel{, }\DIFdelend as decision rules \DIFdelbegin \DIFdel{, }\DIFdelend against production in a common unit.

\begin{table}[t]
\tabsize
\setlength{\tabcolsep}{3.5pt}
\caption{Positioning among methods that support or evaluate design-stage manufacturability decisions}\label{tab:positioning}
\begin{tabular}{@{}L{3.0cm}L{2.5cm}L{3.0cm}L{1.9cm}L{1.6cm}@{}}
\toprule
Approach & Evidence & Output & Scored against production & Abstains \\
\midrule
DFM rules and cost models \citep{boothroyd2010product} & codified process knowledge & guideline compliance, cost & no & no \\
Design margins \citep{eckert2019design} & experience, analysis & buffer to the requirement & rarely & no \\
V\&V, credibility assessment \citep{oberkampf2010verification, asme2018assessing} & physics model, validation tests & accuracy in physical units; required evidence & on test cases & can require tests \\
Bayesian calibration \citep{kennedy2001bayesian} & simulations, experiments & predictive distribution & not as decisions & implicitly \\
Conformal prediction \citep{lei2018distribution} & exchangeable calibration data & interval with coverage & coverage; per decision possible & yes \\
Conformity assessment, OC curves \citep{iso2017decision, bechhofer1954single} & measurement and its uncertainty & accept, reject, undecided; risks & on the measured item & yes \\
Selective classification \citep{chow1970optimum, geifman2017selective} & labelled data & class or reject; risk--coverage & on held-out data & yes \\
ML manufacturability models \citep{ghadai2018learning, yan2026machine} & historical or simulated designs & class or value & on held-out data & rarely \\
This paper & simulations, produced alternatives & verdict; decisive and safe distance in floors by evidence relation & yes, per verdict & yes \\
\botrule
\end{tabular}
\end{table}

%% ----------------------------------------------------------------------
\section{The framework}\label{sec:method}

\subsection{The decision problem and its scope}\label{sec:problem}

A design alternative $a$ is a combination of product geometry and process settings \DIFdelbegin \DIFdel{that is }\DIFdelend produced as a series of parts\DIFdelbegin \DIFdel{. A }\DIFdelend \DIFaddbegin \DIFadd{, over which a }\DIFaddend quality characteristic $c$ varies\DIFdelbegin \DIFdel{over the parts of a series}\DIFdelend . A requirement is an upper limit $R$ on the characteristic, or on its absolute value for characteristics whose target is zero, that a share $p$ of the parts must meet: alternative $a$ is \emph{adequate} for $c$ when the $p$-quantile of its parts satisfies $q_p(a,c)\le R$. Lower or two-sided limits follow by a change of sign or by the absolute deviation from the target\DIFdelbegin \DIFdel{; a tolerance of $\pm t$ on a wall angle about its design angle, for example, is the limit $R=t$ on the absolute deviation}\DIFdelend . At the design stage $q_p$ is unknown. A decision rule $r$ maps the available evidence to an interval $[L_r, U_r]$ for $q_p$ and to a verdict
\begin{equation}
V_r = \begin{cases} \text{meets} & U_r \le R,\\ \text{fails} & L_r > R,\\ \text{trial} & \text{otherwise.} \end{cases}
\label{eq:verdict}
\end{equation}
The three verdicts mirror the decision rules of conformity assessment \citep{iso2017decision}, here for the uncertainty of a design-stage prediction. A rule that returns a point ($L_r=U_r$) always decides. The evidence consists of simulations of $a$ and of the production record of other alternatives, the \emph{calibration alternatives}\DIFdelbegin \DIFdel{; it is used by the process planner or tryout engineer for a further variant or process setting of an existing tool, at tool tryout or before production part approval, and by the product designer for a new family}\DIFdelend . The demonstration decides on $p=0.95$, which for a normal characteristic corresponds to a performance index of about 0.55. Release criteria of series production, a performance index of 1.33 or more \citep{aiag2006production}, concern nonconforming fractions in parts per million that a series of 500 parts cannot estimate without a parametric tail; \DIFdelbegin \DIFdel{the framework reaches them through requirements anchored in capability (}\DIFdelend Section~\ref{sec:res-requirements} \DIFdelbegin \DIFdel{)}\DIFdelend \DIFaddbegin \DIFadd{locates such limits relative to the distances, without deciding them}\DIFaddend .

\subsection{Production floor, effect-to-scatter ratio and minimum resolvable change}\label{sec:floor}

Production evidence comes as series of consecutive parts, each cut into batches of $B$ consecutive parts, and each batch is summarised by a robust centre $m$ (here the 10\,\% trimmed mean). The production floor of characteristic $c$ is
\begin{equation}
F_c = Q_\alpha\bigl\{\,|m_i - m_j| : i< j \text{ batches of the same alternative}\,\bigr\},
\label{eq:floor}
\end{equation}
the $\alpha$-quantile of the difference between two batches of one design. Under a stationary normal model with a standard deviation $\sigma_\mathrm{b}$ between batch means and $\sigma_\mathrm{w}$ between parts, $F_c = 1.96\sqrt{2}\sqrt{\sigma_\mathrm{b}^2+\sigma_\mathrm{w}^2/B}$ at $\alpha=0.95$\DIFdelbegin \DIFdel{: }\DIFdelend \DIFaddbegin \DIFadd{, the form of }\DIFaddend the repeatability limit of precision experiments \DIFdelbegin \DIFdel{\mbox{%DIFAUXCMD
\citep{iso2025repeatability, iso1994accuracy}}\hskip0pt%DIFAUXCMD
, applied to batch centres under within-run conditions}\DIFdelend \DIFaddbegin \DIFadd{\mbox{%DIFAUXCMD
\citep{iso2025repeatability}}\hskip0pt%DIFAUXCMD
. Because it pools batch pairs up to 450 parts apart and grows with their lag (Section~\ref{sec:res-floor}), it is closer to a time-different intermediate-precision limit \mbox{%DIFAUXCMD
\citep{iso2023intermediate} }\hskip0pt%DIFAUXCMD
of batch centres within a run; in the terms of ISO 22514-2 the runs follow a time-dependent model with varying location \mbox{%DIFAUXCMD
\citep{iso2026process}}\hskip0pt%DIFAUXCMD
}\DIFaddend . The reference protocol of this paper is $B=50$, ten batches of a 500-part run, all pairs of batches and $\alpha=0.95$, with per-geometry floors where geometries differ. The floor rests on three assumptions: batches of one design differ only by drift and scatter within a run (A1); the \DIFaddbegin \DIFadd{variation within the }\DIFaddend runs from which it is estimated \DIFdelbegin \DIFdel{are like those of future production }\DIFdelend \DIFaddbegin \DIFadd{is like that within future runs }\DIFaddend (A2); and the measuring system is repeatable and its \DIFaddbegin \DIFadd{measurement }\DIFaddend resolution small against the floor (A3, checked in Section~\ref{sec:res-floor}).

Four properties matter for its use. It has the unit of the characteristic, so distances in floors are on a common scale for all characteristics and predictors, \DIFdelbegin \DIFdel{although }\DIFdelend \DIFaddbegin \DIFadd{though }\DIFaddend equal distances do not mean equal nonconforming fractions \DIFdelbegin \DIFdel{, because the ratio of the floor to the scatter of the parts differs between characteristics }\DIFdelend (Table~\ref{tab:floors}). It is a property of a \DIFdelbegin \DIFdel{production }\DIFdelend protocol: under drift it grows with the time between the batches compared \DIFdelbegin \DIFdel{, and it }\DIFdelend \DIFaddbegin \DIFadd{and }\DIFaddend shrinks with $B$\DIFdelbegin \DIFdel{; }\DIFdelend \DIFaddbegin \DIFadd{, and }\DIFaddend $\sigma_\mathrm{b}$ and $\sigma_\mathrm{w}$ are reported so that it can be recomputed\DIFdelbegin \DIFdel{for another protocol (Table~\ref{tab:floors})}\DIFdelend . It contains whatever drifts between batches, \DIFdelbegin \DIFdel{the process and the }\DIFdelend \DIFaddbegin \DIFadd{process and }\DIFaddend measuring system alike, so without a reference artefact scanned through the series it is a production-and-measurement floor. And it excludes the variation between runs, which makes it a lower bound of the variation of series production. \DIFdelbegin \DIFdel{For a }\DIFdelend \DIFaddbegin \DIFadd{A }\DIFaddend new family, whose floor is unknown before it is produced, \DIFaddbegin \DIFadd{is scored in }\DIFaddend the floor of the produced family\DIFdelbegin \DIFdel{is used. Because }\DIFdelend \DIFaddbegin \DIFadd{, and because }\DIFaddend the floor resolves batch centres\DIFdelbegin \DIFdel{while the requirement concerns a quantile}\DIFdelend , the reproducibility of the \DIFaddbegin \DIFadd{decided }\DIFaddend quantile is reported alongside.

Two alternatives $a$ and $b$ are resolvable for $c$ when their effect-to-scatter ratio
\begin{equation}
\mathrm{ESR}_c(a,b) = |\bar m_a - \bar m_b| \,/\, F_c
\label{eq:esr}
\end{equation}
is at least one, $\bar m$ being the centre of all parts of an alternative; this is a practical threshold, not a test. For a continuous design or process parameter $x$ with a local sensitivity $s_c = \partial \bar m/\partial x$, the minimum resolvable change is $F_c/|s_c|$: a change of $x$ by less moves the centre of the characteristic by less than two batches of one design differ. It is local \DIFdelbegin \DIFdel{, a statement about centres, and }\DIFdelend \DIFaddbegin \DIFadd{and concerns centres; }\DIFaddend a change below it can still move the tail\DIFdelbegin \DIFdel{of the distribution}\DIFdelend .

\subsection{Decision rules}\label{sec:rules}

\DIFdelbegin \DIFdel{Thirteen }\DIFdelend \DIFaddbegin \DIFadd{Fourteen }\DIFaddend rules span what a design team can do with its evidence (Table~\ref{tab:rules}); \DIFdelbegin \DIFdel{each is defined for any set of calibration alternatives, and }\DIFdelend the first seven \DIFdelbegin \DIFdel{carry the results in the main text. The nominal simulation }\DIFdelend \DIFaddbegin \DIFadd{and the force trend M1sn carry the main results. Beside the nominal simulation $s_0$ }\DIFaddend (M0)\DIFdelbegin \DIFdel{uses the simulation alone. The }\DIFdelend \DIFaddbegin \DIFadd{, the }\DIFaddend bias-corrected simulation (M1) adds the median offset between the calibration alternatives' $q_p$ and their \DIFdelbegin \DIFdel{nominal simulation}\DIFdelend \DIFaddbegin \DIFadd{$s_0$}\DIFaddend ; the scaled simulation (M1s) fits $q_p = a + b\,s_0$ to them by least squares, a linear multi-fidelity correction \citep{kennedy2000predicting} that rescales a simulated trend of the wrong size. The envelope rule (M2) adds the median offset and the largest absolute residual of the calibration alternatives to the upper end of the simulation envelope over plausible incoming conditions\DIFaddbegin \DIFadd{, which vary sheet thickness and friction but not material properties}\DIFaddend . M5 models the offset between $q_p$ and the envelope as a Gaussian process \citep{rasmussen2006gaussian} with a 90\,\% predictive interval, whose noise variance is bounded below by the variance between the \DIFdelbegin \DIFdel{produced }\DIFdelend \DIFaddbegin \DIFadd{calibration alternatives' }\DIFaddend series of one geometry and force, so that a new series of a produced setting is not predicted more precisely than produced series reproduce each other; M5n is the same process of $q_p$ without the simulation. The nearest produced setting (NN) is a capability-database look-up that uses neither simulation nor learning. The \DIFdelbegin \DIFdel{six }\DIFdelend \DIFaddbegin \DIFadd{seven }\DIFaddend further rules include a friction-calibrated simulation \citep{kennedy2001bayesian} (Mc)\DIFdelbegin \DIFdel{and }\DIFdelend \DIFaddbegin \DIFadd{, }\DIFaddend part-level gradient boosting \citep{friedman2001greedy} with a jackknife+ interval \citep{barber2021predictive}, with and without the simulation (M3, M3n)\DIFaddbegin \DIFadd{, and the least-squares line of $q_p$ in force (M1sn)}\DIFaddend . The nominal simulation and the envelope do not depend on the lubrication pattern, so \DIFdelbegin \DIFdel{a lubrication variant of a produced setting }\DIFdelend \DIFaddbegin \DIFadd{within a geometry they are functions of the force alone: M1sn is the counterpart of M1s without the simulation, and a lubrication variant }\DIFaddend gives M0, M1, M1s, M2 and M5 no simulated information that distinguishes it from its siblings.

\DIFdelbegin \DIFdel{The margins are labelled by what they guarantee. }\DIFdelend With $n$ calibration alternatives, the largest-residual margin covers a new exchangeable residual with probability $n/(n+1)$ for a fixed offset and at least $(n-1)/(n+1)$ for an estimated one\DIFdelbegin \DIFdel{, because the new residual is then covered whenever it falls within the range of the calibration residuals}\DIFdelend ; it is the split-conformal quantile \citep{lei2018distribution} at the highest attainable level, not a 90\,\% interval. \DIFdelbegin \DIFdel{In general, a }\DIFdelend \DIFaddbegin \DIFadd{A }\DIFaddend distribution-free interval cannot certify a miss rate below $1/(n+1)$, so a miss rate of 5\,\% needs at least 19 produced alternatives in the relevant relation\DIFdelbegin \DIFdel{and one of 10\,\% needs nine. Under exchangeability }\DIFdelend \DIFaddbegin \DIFadd{. With six to nine calibration alternatives }\DIFaddend the jackknife+ \DIFdelbegin \DIFdel{interval covers with probability at least 0.8 }\DIFdelend \DIFaddbegin \DIFadd{bounds }\DIFaddend at the level of 0.9 \DIFdelbegin \DIFdel{used here \mbox{%DIFAUXCMD
\citep{barber2021predictive}}\hskip0pt%DIFAUXCMD
, and the }\DIFdelend \DIFaddbegin \DIFadd{fall at or beyond the extreme leave-one-out values, which are used instead; under exchangeability these guarantee a coverage of $1-2/(n+1)$, 0.71 to 0.80 \mbox{%DIFAUXCMD
\citep{barber2021predictive}}\hskip0pt%DIFAUXCMD
, approximately because the residuals are centred on their median. The }\DIFaddend Gaussian-process interval is model-based. The level \DIFdelbegin \DIFdel{is a decision variable that }\DIFdelend trades decisiveness for safety, as the cost of errors would set the threshold of a reject option \citep{chow1970optimum}; 0.9 is used here, \DIFdelbegin \DIFdel{and }\DIFdelend 0.8 and 0.95 \DIFdelbegin \DIFdel{are reported }\DIFdelend in the robustness analysis.

\begin{table}[t]
\tabsize
\setlength{\tabcolsep}{3.5pt}
\DIFdelbeginFL %DIFDELCMD < \caption{The decision rules; the first seven carry the results in the main text}%%%
\DIFdelendFL \DIFaddbeginFL \caption{The decision rules; the first seven and M1sn carry the main results}\DIFaddendFL \label{tab:rules}
\DIFdelbeginFL %DIFDELCMD < \begin{tabular}{@{}L{3.3cm}L{5.0cm}L{2.0cm}L{1.6cm}@{}}
%DIFDELCMD < \toprule
%DIFDELCMD < Rule & Prediction of $q_p$ & Interval & Production evidence \\
%DIFDELCMD < \midrule
%DIFDELCMD < M0 nominal simulation & simulation at nominal sheet thickness and friction & point & none \\
%DIFDELCMD < M1 bias-corrected simulation & M0 plus the median offset & point & $q_p$ \\
%DIFDELCMD < M1s scaled simulation & $a+b\,$M0, least squares on the calibration alternatives & point & $q_p$ \\
%DIFDELCMD < M2 envelope + residual margin & upper end of the envelope plus the median offset & largest residual & $q_p$ \\
%DIFDELCMD < M5 GP correction of the envelope & envelope plus a Gaussian-process offset over the descriptors & GP, 90\,\% & $q_p$ \\
%DIFDELCMD < M5n GP of $q_p$ & Gaussian process of $q_p$ over the descriptors & GP, 90\,\% & $q_p$ \\
%DIFDELCMD < NN nearest produced setting & mean $q_p$ at the nearest produced process setting of the family & point & $q_p$ \\
%DIFDELCMD < \midrule
%DIFDELCMD < M0w process-window worst case & lowest to highest simulated value over the incoming conditions & range & none \\
%DIFDELCMD < Mc friction-calibrated simulation & simulation with friction calibrated on the draw-in, plus the median offset & point & centres, $q_p$ \\
%DIFDELCMD < M3 simulation + boosting & matched simulation plus a boosting model of the discrepancy, part level & jackknife+ & parts \\
%DIFDELCMD < M3n boosting, no simulation & boosting model of the characteristic, part level & jackknife+ & parts \\
%DIFDELCMD < M1n median of the produced & median $q_p$ of the calibration alternatives & point & $q_p$ \\
%DIFDELCMD < M2n median + residual margin & M1n with the largest-residual margin & largest residual & $q_p$ \\
%DIFDELCMD < \botrule
%DIFDELCMD < \end{tabular}
%DIFDELCMD < %%%
\DIFdelendFL \DIFaddbeginFL \begin{tabular}{@{}L{3.5cm}L{4.9cm}L{1.9cm}L{1.4cm}@{}}
\toprule
Rule & Prediction of $q_p$ & Interval & Evidence \\
\midrule
M0 nominal simulation & simulation at nominal thickness and friction & point & none \\
M1 bias-corrected simulation & M0 plus the median offset & point & $q_p$ \\
M1s scaled simulation & $a+b\,$M0, least squares & point & $q_p$ \\
M2 envelope + residual margin & envelope top plus the median offset & largest residual & $q_p$ \\
M5 GP-corrected envelope & envelope plus a GP offset over the descriptors & GP, 90\,\% & $q_p$ \\
M5n GP of $q_p$ & GP of $q_p$ over the descriptors & GP, 90\,\% & $q_p$ \\
NN nearest produced setting & mean $q_p$ at the nearest produced setting & point & $q_p$ \\
\midrule
M0w process-window worst case & simulated range over the incoming conditions & range & none \\
Mc friction-calibrated simulation & friction calibrated on the draw-in, plus offset & point & centres, $q_p$ \\
M3 simulation + boosting & matched simulation plus boosted discrepancy & jackknife+ & parts \\
M3n boosting, no simulation & boosted model of the characteristic & jackknife+ & parts \\
M1n median of the produced & median $q_p$ of the calibration alternatives & point & $q_p$ \\
M1sn force trend & $a+b\,$force, least squares & point & $q_p$ \\
M2n median + residual margin & M1n with the largest-residual margin & largest residual & $q_p$ \\
\botrule
\end{tabular}
\DIFaddendFL \footnotetext{\DIFaddbeginFL \DIFaddFL{Fitted on the calibration alternatives. }\DIFaddendFL Descriptors: geometry indicator, blank-holder force and lubrication rank (the order of the oil amount). GP: \DIFdelbeginFL \DIFdelFL{noise variance bounded below by the variance between the series of one geometry and force}\DIFdelendFL \DIFaddbeginFL \DIFaddFL{Gaussian process}\DIFaddendFL . \DIFaddbeginFL \DIFaddFL{M3 and M3n at part level; }\DIFaddendFL M3n is named M4 in the released code\DIFdelbeginFL \DIFdelFL{; }\DIFdelendFL \DIFaddbeginFL \DIFaddFL{. }\DIFaddendFL M6, a conformalised GP, \DIFdelbeginFL \DIFdelFL{appears only in the }\DIFdelendFL \DIFaddbeginFL \DIFaddFL{is a }\DIFaddendFL robustness \DIFdelbeginFL \DIFdelFL{analysis}\DIFdelendFL \DIFaddbeginFL \DIFaddFL{variant (Table~\ref{tab:distances})}\DIFaddendFL .}
\end{table}

\subsection{Operating characteristics: decisive and safe distances}\label{sec:distances}

A rule is scored by judging each produced alternative as a new design at requirements placed at a known distance from the truth,
\begin{equation}
R = q_p(a,c) + d\,F_c ,
\label{eq:requirement}
\end{equation}
so that $a$ is adequate exactly when $d\ge 0$. For a case $(a,c)$ with interval $[L,U]$, let $u=(U-q_p)/F_c$ and $v=(q_p-L)/F_c$ be the distances by which the interval excludes the truth from above and below. At $d=+\delta$ the verdict is correct when $u\le\delta$ and a false reject when $-v>\delta$; at $d=-\delta$ it is correct when $v<\delta$ and a false accept when $-u\ge\delta$. Let $\kappa(\delta)$ and $\omega(\delta)$ be the shares of correct and of wrong verdicts among all feasible verdicts at $|d|=\delta$: every case gives one verdict on an adequate alternative and, where $q_p-\delta F_c\ge0$, one on a failing alternative, and every verdict has equal weight. As $\delta$ grows, failing-side requirements become infeasible for characteristics close to zero, so the failing side carries half the weight up to a few floors and a third at about 100 floors (98 of 126 failing-side requirements remain feasible at 16 floors and 63 at 108). The \emph{decisive distance} $\delta_\mathrm{d}$ of a rule is the smallest $\delta$ beyond which $\kappa\ge 0.95$, and its \emph{safe distance} $\delta_\mathrm{s}$ the smallest $\delta$ beyond which $\omega\le 0.05$.

The distances are exact \DIFaddbegin \DIFadd{up to the 0.05-floor grid on which they are evaluated}\DIFaddend : $\kappa$ and $\omega$ change only where $\delta$ crosses one of the case's $u$, $v$, $-u$ or $-v$, so the distances are order statistics of these exclusion distances; for a rule that always decides, $\delta_\mathrm{d}=\delta_\mathrm{s}$ is close to the 90\,\% quantile of its absolute error in floors. Always $\delta_\mathrm{s}\le\delta_\mathrm{d}$, and between the two the rule sends designs to trial rather than misleading; one-sided distances for failing alternatives (false accepts) and adequate ones (false rejects) matter when the errors cost differently. The distances are points on operating characteristics over a population of design situations: $\delta_\mathrm{d}$ is the analogue of the indifference zone of a selection procedure \citep{bechhofer1954single} at a probability of correct selection of 0.95\DIFdelbegin \DIFdel{, and the pair plays the role of the acceptable and rejectable quality levels of an acceptance plan}\DIFdelend . A trial is \DIFdelbegin \DIFdel{itself }\DIFdelend not error-free: a short run reproduces the 95th percentile only to about one floor \DIFaddbegin \DIFadd{within a run }\DIFaddend (Section~\ref{sec:res-floor}) \DIFaddbegin \DIFadd{and, being a run, also carries the variation between runs}\DIFaddend , so no decision \DIFdelbegin \DIFdel{can be relied on }\DIFdelend \DIFaddbegin \DIFadd{is reliable }\DIFaddend closer to the requirement than the trial \DIFdelbegin \DIFdel{that would replace }\DIFdelend \DIFaddbegin \DIFadd{replacing }\DIFaddend it.

\subsection{Evidence relation, calibration budget and applicability guard}\label{sec:budget}

How a new design relates to the produced alternatives is part of the decision situation. With a \emph{sibling}, an alternative at the same process setting has been produced \DIFdelbegin \DIFdel{with another variation that }\DIFdelend \DIFaddbegin \DIFadd{that differs only in the factor varied at a fixed setting (here the lubrication pattern), which }\DIFaddend production may not resolve; a \emph{new process setting} lies between produced settings (interpolation) or outside them (extrapolation); a \emph{new family} shares no produced alternative. \DIFdelbegin \DIFdel{Beyond a factorial, the relation can be graded by a distance in descriptor space or by the effect-to-scatter ratio between the predicted characteristic and the nearest produced alternative (Section~\ref{sec:reuse}). }\DIFdelend The calibration budget is the dependence of the distances on the number $k$ of produced alternatives, over all subsets of size $k$ classified by relation. An applicability guard returns \emph{trial} when the new design lies outside the calibration: its geometry is not calibrated (family guard), a descriptor lies outside the calibrated range (range guard), or its simulated $q_p$ lies outside the calibration alternatives' simulated values (novelty guard), in the spirit of a valid input domain \citep{malak2010using}.

\subsection{Procedure}\label{sec:procedure}

Figure~\ref{fig:procedure} turns the framework into six steps: steps 1 to 5 once per design family, repeated when production evidence is added, and step 6 for every new design. \DIFdelbegin \DIFdel{Step }\DIFdelend \DIFaddbegin \DIFadd{Steps }\DIFaddend 1 \DIFdelbegin \DIFdel{frames }\DIFdelend \DIFaddbegin \DIFadd{to 4 frame }\DIFaddend the decision (family, alternatives, characteristics, form of the requirement\DIFdelbegin \DIFdel{and }\DIFdelend \DIFaddbegin \DIFadd{, }\DIFaddend conformance level)\DIFdelbegin \DIFdel{; step 2 measures }\DIFdelend \DIFaddbegin \DIFadd{, measure }\DIFaddend the floor with a stated protocol\DIFdelbegin \DIFdel{; step 3 states which differences between alternatives, and which local changes of a parameter, production can resolve; and step 4 calibrates }\DIFdelend \DIFaddbegin \DIFadd{, state which differences production resolves and calibrate }\DIFaddend the candidate rules with and without the simulation. Step 5 qualifies them by their distances for each evidence relation and fixes for each rule and relation a \emph{guard band} $g$: the smallest observable margin between the predicted interval and the requirement beyond which at least 95\,\% of the decided verdicts were correct\DIFaddbegin \DIFadd{. Being conditional on what is observed, $g$ depends on the distribution of requirement distances that the team expects, not on rule and relation alone. The reference prior weights the 47 distances of the evaluation grid within $\pm20$ floors equally (0}\DIFaddend , \DIFdelbegin \DIFdel{for requirements within 20 floors of the truth}\DIFdelend \DIFaddbegin \DIFadd{$\pm0.5$ to $\pm10$ in steps of 0.5, $\pm12$, $\pm15$ and $\pm20$; 87\,\% of the weight within $\pm10$), drops infeasible requirements and searches $g$ in steps of 0.25 floors up to 20; it represents requirements set from a few to some tens of floors from the true 95th percentile, denser where verdicts are hardest}\DIFaddend . If the distances are too large for the requirements of coming designs, another variant is produced. Step 6 establishes the relation of the new design, applies the rule qualified for it and the guards, widens the rule's interval by $g$ floors and returns meets or fails with the margin in floors, or a trial when a guard fires, no rule is qualified or the requirement lies inside the widened interval\DIFaddbegin \DIFadd{; Section~\ref{sec:res-procedure} reports $g$ under other priors and works one decision through}\DIFaddend .

\begin{figure}[t]
\centering
\input{fig_procedure}
\caption{The procedure: steps 1--5 once per design family, step 6 for every new design; dashed boxes are the evidence it uses}\label{fig:procedure}
\end{figure}

%% ----------------------------------------------------------------------
\section{Demonstration on deep drawing}\label{sec:demo}

\subsection{Alternatives and production data}\label{sec:data}

The production data are the Real Deep Drawing and Cutting dataset (RDDAC) of the University of Stuttgart \citep{baum2026rddac, baum2026statistical}. Square blanks of DP600 steel, 210\,mm wide and about 1\,mm thick, were deep drawn into cups (OP10) and cut (OP20) on one tool set with two geometries, concave and convex side walls, which differ in base curvature, wall angle and bottom radius at once. Each geometry was run at three blank-holder forces (100, 300 and 500\,kN) and three lubrication patterns (coarse, medium and fine, with series medians of 1.15--1.48, 0.92--1.23 and 0.85--1.13\,g/m$^2$ of oil): 18 alternatives, each produced as one series of 500 consecutive parts, the change of geometry being the only change of product design. Every part has press signals, sheet and oil-film traverses and 3D scans after both operations\DIFdelbegin \DIFdel{; 10 of the 9,000 parts carry no scan}\DIFdelend . The series warm up by 1.0--5.0\,K over their first 150 parts from punch temperatures of 20.5--29.6\,$^\circ$C, their median sheet thickness \DIFdelbegin \DIFdel{ranges from }\DIFdelend \DIFaddbegin \DIFadd{is }\DIFaddend 0.986\DIFdelbegin \DIFdel{to 0.993}\DIFdelend \DIFaddbegin \DIFadd{--0.993}\DIFaddend \,mm, and the 100\,kN series ran at 47\,mm/s against 78--79\,mm/s. The documentation states neither the order and dates of the series nor whether the parts were scanned in production order.

\subsection{Simulations}\label{sec:sims}

The simulations come from the Deep Drawing and Cutting Simulations dataset (DDACS) of the same group \citep{baum2025ddacs, heinzelmann2025comprehensive}: LS-DYNA quarter models with nodes about 1\,mm apart, a scaled DP600 flow curve, constant Coulomb friction and springback after both operations; the yield criterion and unloading model, which govern springback \citep{yoshida2002model}, are not documented, and the stroke speed is not represented. For the RDDAC tools, 396 simulations vary the sheet thickness (0.95--1.00\,mm) and the friction coefficient (0.05--0.15) at each geometry and force. A part is matched to the simulation with the nearest thickness and the friction coefficient that the dataset's linear mapping assigns to its oil film, an assumption of the dataset's authors on which Mc, M3 and every matched simulation depend. The nominal simulation of an alternative is at 0.99\,mm, the grid value nearest the median thickness, and friction 0.10.

\subsection{Quality characteristics and their measurement}\label{sec:qcs}

Seven characteristics were extracted with the same definitions from the scans and the simulated nodes, referred to the plane of the cup bottom: the mid-side and corner draw-in and the flange waviness, process indicators on the flange that the cut removes, and four characteristics a drawing would carry: the wall angle (design angle 20$^\circ$ concave, 10$^\circ$ convex), the cup depth, and the arm angle and bottom dome of the cut part, the last two with requirements on absolute values. The four sides of a part measure several of them twice, which exposes the measurement: the flange width along the scan lines scatters 4 to 18 times more than across them, so the mid-side draw-in is taken across the scan lines, and the wall angle is the mean of the east and west walls\DIFdelbegin \DIFdel{(the north wall carries uncorrelated scatter, the south wall is shadowed)}\DIFdelend . The mid-side draw-in is quantised in steps of 0.019\,mm (0.38 floors) by the pixel pitch. Differences between scans and simulations, such as the scanned surface against the shell mid-surface, enter the offsets between the two but not the floors.

\subsection{Evaluation design}\label{sec:evaldesign}

Every produced alternative is judged as a new design under grouped cross-validation \citep{roberts2017cross}, one scheme per evidence relation. \emph{New variant}: the alternative alone is held out and the other eight of its geometry calibrate, two siblings at its force among them. \emph{New process setting}: the three alternatives of one force of a geometry are held out and the six others calibrate; 300\,kN is interpolated, 100 and 500\,kN extrapolated. \emph{New family}: one geometry calibrates the rules for the other (two transfers). Pooled versions of the first two add the other geometry's alternatives to the calibration set. Every calibrated rule is evaluated with and without the simulation (M1/M1n, \DIFaddbegin \DIFadd{M1s/M1sn, }\DIFaddend M2/M2n, M3/M3n, M5/M5n), \DIFdelbegin \DIFdel{so that the paired difference isolates }\DIFdelend \DIFaddbegin \DIFadd{isolating }\DIFaddend what the simulation contributes, and the nearest produced setting is the baseline without simulation and learning. A scope comprises 126 cases (18 alternatives, 7 characteristics).

The \DIFdelbegin \DIFdel{distances are computed exactly. Their intervals }\DIFdelend \DIFaddbegin \DIFadd{intervals of the distances }\DIFaddend come from 1000 resamples of the held-out alternatives within each geometry, each with a block-bootstrap replicate of every true $q_p$, common to all rules so that differences get paired intervals; they hold the fitted rules fixed. A second bootstrap \DIFdelbegin \DIFdel{resamples whole lubrication patterns within each geometry }\DIFdelend \DIFaddbegin \DIFadd{refits every alternative-level rule after drawing the three lubrication patterns of each geometry with replacement, at least two of them distinct}\DIFaddend , which keeps every case in its evidence relation \DIFdelbegin \DIFdel{, and refits every alternative-level rule }\DIFdelend (Table~\ref{tab:robust}). \DIFaddbegin \DIFadd{Its 200 resamples contain 48 of the 49 distinct configurations; in three of four draws a pattern is duplicated, so that a new variant keeps one distinct sibling and five distinct calibration alternatives, and the duplicated series enter the noise floor of the Gaussian processes as identical pairs. Its intervals thus show the effect of losing a pattern, not the sampling uncertainty of the reported values.
}\DIFaddend 

\subsection{Computation and development record}\label{sec:computation}

The constants of the floor and the requirement grid were fixed before the decision results were seen, as the repository history records. The rules were not: M0--M5 were defined first; M0w, Mc, M1n, NN, M2n, M5n, the grouped schemes, the exact estimator and the noise floor of the Gaussian processes were added after the first results; \DIFdelbegin \DIFdel{and }\DIFdelend a review of that version added M1s, the jackknife+ intervals, the between-run noise floor and the analyses of Sections~\ref{sec:res-drivers} and \ref{sec:res-procedure}\DIFaddbegin \DIFadd{; and a second review added M1sn, the guard bands under other priors and the new family in the produced family's floor}\DIFaddend . The comparison of rules is therefore exploratory and carries the optimism of a selection made on the evaluation data \citep{cawley2010over}; every variant tried is reported in the paper or \DIFaddbegin \DIFadd{in }\DIFaddend the archived results \DIFdelbegin \DIFdel{. The pipeline reproduces every result from the cached feature tables in about four hours on four cores}\DIFdelend \DIFaddbegin \DIFadd{\mbox{%DIFAUXCMD
\citep{canseven2026archive}}\hskip0pt%DIFAUXCMD
, which Table~\ref{tab:mapping} maps to the text}\DIFaddend .

%% ----------------------------------------------------------------------
\section{Results}\label{sec:results}

\subsection{What the production floor contains}\label{sec:res-floor}

In 95\,\% of batch pairs, two batches of 50 consecutive parts of one alternative differ by up to 0.051\,mm in mid-side draw-in and 0.078$^\circ$ in wall angle (Table~\ref{tab:floors}); one floor is \DIFdelbegin \DIFdel{0.66 to 2.2 }\DIFdelend \DIFaddbegin \DIFadd{0.7 to 1.7 }\DIFaddend long-term standard deviations of the parts. The floors are set by drift within the run: without the time order they shrink by a factor of 2.1 to 4.5, \DIFaddbegin \DIFadd{and }\DIFaddend batches one position apart differ by 0.51--0.80 floors and nine apart by 1.1--1.5 \DIFdelbegin \DIFdel{, and the }\DIFdelend \DIFaddbegin \DIFadd{(averages over the two geometries). The }\DIFaddend normal model of Section~\ref{sec:floor} reproduces the floors to within 12\,\%\DIFaddbegin \DIFadd{; for subgroups of five parts, as sampled inspection would measure, it gives floors 1.1 to 2.1 times as large (median 1.3)}\DIFaddend . Independent part-to-part noise, scanner repeatability included, would produce 15--33\,\% of each floor, and the \DIFaddbegin \DIFadd{measurement }\DIFaddend resolution is small against the floor except for the mid-side draw-in; \DIFdelbegin \DIFdel{assumption }\DIFdelend A3 holds for repeatability and resolution, but the stability of the scanner cannot be checked without a reference artefact. Nor do the process signals identify the source of the drift: batch centres of press force, punch temperature, sheet thickness and oil film predict it out of fold with $R^2$ of $-0.23$ to $0.20$\DIFdelbegin \DIFdel{, and removing the predicted drift changes the floors by $-16$ to $+12$\,\%}\DIFdelend . Variation between runs is not in the floor: the three lubrication series of one geometry and force differ by up to 1.2 (waviness) to 9.8 floors (dome), and a floor taken across them is 1.5 to 9.1 times the reference. \DIFdelbegin \DIFdel{The }\DIFdelend \DIFaddbegin \DIFadd{Within a run the }\DIFaddend decided statistic reproduces to about one floor \DIFdelbegin \DIFdel{: }\DIFdelend \DIFaddbegin \DIFadd{(}\DIFaddend the 95th percentiles of two batches of 100 parts differ by 0.92--1.49 floors in 95\,\% of pairs\DIFaddbegin \DIFadd{)}\DIFaddend , so distances below about one floor are beyond what production \DIFdelbegin \DIFdel{, or a trial , can confirm}\DIFdelend \DIFaddbegin \DIFadd{can confirm, and a trial run, which also carries the variation between runs, confirms less still}\DIFaddend .

%% BEGIN GENERATED floors
\begin{table}[t]
\tabsize
\setlength{\tabcolsep}{2.0pt}
\caption{Production floor per characteristic, pooled and per geometry, with its components}\label{tab:floors}
\DIFdelbeginFL %DIFDELCMD < \begin{tabular}{@{}lrrrrrrrr@{}}
%DIFDELCMD < \toprule
%DIFDELCMD < Characteristic & $F$ & 95\,\% interval & Concave & Convex & $F/\sigma_\mathrm{LT}$ & $\sigma_\mathrm{b}/F$ & $\sigma_\mathrm{w}/F$ & Drift \\
%DIFDELCMD < \midrule
%DIFDELCMD < Draw-in, mid-side (mm) & 0.051 & 0.045--0.061 & 0.050 & 0.052 & 1.2/1.5 & 0.32/0.34 & 0.91/0.66 & 3.3 \\
%DIFDELCMD < Draw-in, corner (mm) & 0.046 & 0.035--0.052 & 0.053 & 0.030 & 1.7/1.3 & 0.31/0.35 & 0.65/0.77 & 4.5 \\
%DIFDELCMD < Flange waviness (mm) & 0.0028 & 0.0024--0.0031 & 0.0037 & 0.0021 & 2.1/0.8 & 0.33/0.30 & 0.55/1.31 & 2.5 \\
%DIFDELCMD < Wall angle ($^\circ$) & 0.078 & 0.066--0.10 & 0.081 & 0.078 & 1.6/0.7 & 0.32/0.26 & 0.68/1.39 & 2.1 \\
%DIFDELCMD < Arm angle, cut ($^\circ$) & 0.11 & 0.094--0.12 & 0.12 & 0.100 & 1.9/0.9 & 0.36/0.31 & 0.46/1.22 & 2.6 \\
%DIFDELCMD < Cup depth (mm) & 0.011 & 0.0096--0.012 & 0.0087 & 0.012 & 2.2/1.7 & 0.31/0.34 & 0.65/0.52 & 3.9 \\
%DIFDELCMD < Bottom dome, cut (mm) & 0.0070 & 0.0061--0.0081 & 0.0034 & 0.0082 & 1.4/1.5 & 0.33/0.36 & 0.76/0.55 & 3.9 \\
%DIFDELCMD < \botrule
%DIFDELCMD < \end{tabular}
%DIFDELCMD < %%%
\DIFdelendFL \DIFaddbeginFL \begin{tabular}{@{}lrrrrrrrr@{}}
\toprule
Characteristic & $F$ & 95\,\% interval & Concave & Convex & $F/\sigma_\mathrm{LT}$ & $\sigma_\mathrm{b}/F$ & $\sigma_\mathrm{w}/F$ & Drift \\
\midrule
Draw-in, mid-side (mm) & 0.051 & 0.045--0.061 & 0.050 & 0.052 & 1.0/1.4 & 0.32/0.34 & 0.91/0.66 & 3.3 \\
Draw-in, corner (mm) & 0.046 & 0.035--0.052 & 0.053 & 0.030 & 1.4/1.2 & 0.31/0.35 & 0.65/0.77 & 4.5 \\
Flange waviness (mm) & 0.0028 & 0.0024--0.0031 & 0.0037 & 0.0021 & 1.6/0.7 & 0.33/0.30 & 0.55/1.31 & 2.5 \\
Wall angle ($^\circ$) & 0.078 & 0.066--0.10 & 0.081 & 0.078 & 1.3/0.7 & 0.32/0.26 & 0.68/1.39 & 2.1 \\
Arm angle, cut ($^\circ$) & 0.11 & 0.094--0.12 & 0.12 & 0.100 & 1.7/0.8 & 0.36/0.31 & 0.46/1.22 & 2.6 \\
Cup depth (mm) & 0.011 & 0.0096--0.012 & 0.0087 & 0.012 & 1.4/1.6 & 0.31/0.34 & 0.65/0.52 & 3.9 \\
Bottom dome, cut (mm) & 0.0070 & 0.0061--0.0081 & 0.0034 & 0.0082 & 1.2/1.6 & 0.33/0.36 & 0.76/0.55 & 3.9 \\
\botrule
\end{tabular}
\DIFaddendFL \footnotetext{Batches of 50 consecutive parts, 95\,\% quantile of the difference of \DIFdelbeginFL \DIFdelFL{batch centres}\DIFdelendFL \DIFaddbeginFL \DIFaddFL{their 10\,\% trimmed means}\DIFaddendFL ; interval from 2000 \DIFaddbeginFL \DIFaddFL{batch }\DIFaddendFL resamples\DIFdelbeginFL \DIFdelFL{of batches within series}\DIFdelendFL . Ratios per geometry (concave/convex)\DIFaddbeginFL \DIFaddFL{, classical statistics pooled over the series as root mean squares}\DIFaddendFL : $\sigma_\mathrm{LT}$, \DIFdelbeginFL \DIFdelFL{median long-term }\DIFdelendFL standard deviation of the parts \DIFaddbeginFL \DIFaddFL{of a series}\DIFaddendFL ; $\sigma_\mathrm{b}$, \DIFdelbeginFL \DIFdelFL{standard deviation }\DIFdelendFL between batch means\DIFaddbeginFL \DIFaddFL{, }\DIFaddendFL corrected for $\sigma_\mathrm{w}^2/50$; $\sigma_\mathrm{w}$, within batches. Drift: floor over \DIFdelbeginFL \DIFdelFL{the floor }\DIFdelendFL \DIFaddbeginFL \DIFaddFL{that }\DIFaddendFL of randomly permuted series.}
\end{table}
%% END GENERATED floors

\subsection{What production resolves and how far the simulation misses it}\label{sec:res-sim}

All 63 geometry contrasts exceed the floor and blank-holder force is resolved in 105 of 126 contrasts, but lubrication only in 51 of 126 (42 with a bootstrap probability of at least 0.95), so a lubrication variant is often a near-replicate of its siblings. The minimum resolvable force change is 21--67\,kN between 100 and 300\,kN for most characteristics and lies beyond either force interval for the wall angle and the dome. The nominal simulation misses production by far more than a floor: its median offset reaches 109 floors, its force effect on the draw-in is 2.7 to 5.1 times the measured one, and the cut arms of the concave cups bend downwards in 89--100\,\% of the simulations but upwards in every part. Part of the offset is definitional: the cup depth differs by a median 0.90 and 0.95\,mm, consistent with the scanned surface against the shell mid-surface, and removing this offset alone brings the decisive distance of the nominal simulation from 108 to 65 floors.

\subsection{Decisions by evidence relation}\label{sec:res-decisions}

Table~\ref{tab:distances} and Fig.~\ref{fig:decisions} score the rules. The nominal simulation \DIFdelbegin \DIFdel{needs requirements }\DIFdelend \DIFaddbegin \DIFadd{decides at }\DIFaddend 108 floors \DIFdelbegin \DIFdel{from the truth (bootstrap interval }\DIFdelend \DIFaddbegin \DIFadd{(}\DIFaddend 101--115)\DIFdelbegin \DIFdel{before 95\,\% of its verdicts are right; it errs }\DIFdelend \DIFaddbegin \DIFadd{, erring }\DIFaddend mostly by rejecting adequate alternatives\DIFdelbegin \DIFdel{(safe distance 17 floors for false accepts, 109 for false rejects)}\DIFdelend .

\emph{A new variant with a produced sibling.} \DIFdelbegin \DIFdel{With the other eight alternatives of the family produced, every }\DIFdelend \DIFaddbegin \DIFadd{Every }\DIFaddend calibrated rule improves on \DIFdelbegin \DIFdel{the nominal simulation}\DIFdelend \DIFaddbegin \DIFadd{it}\DIFaddend , and the simplest is the most decisive: the nearest produced setting, the mean $q_{95}$ of the two siblings\DIFdelbegin \DIFdel{at the same force}\DIFdelend , is right beyond 3.4 floors (2.0--4.5), 0.17\,mm of mid-side draw-in\DIFdelbegin \DIFdel{. The scaled simulation follows at 4.9 floors}\DIFdelend ; the Gaussian processes decide at 7.2 floors \DIFdelbegin \DIFdel{and }\DIFdelend \DIFaddbegin \DIFadd{but }\DIFaddend are safe below one floor\DIFdelbegin \DIFdel{; the envelope rule is safe at 0.15 floors but sends a third of the designs to trial at 10 floors and decides only at 23. Bias-correcting the nominal simulation (M1, 15 floors) is worse than the median of the produced alternatives without it (M1n, 9.1)}\DIFdelend .

\emph{A new process setting.} With every alternative at the new design's force held out, all rules lose\DIFdelbegin \DIFdel{. The }\DIFdelend \DIFaddbegin \DIFadd{: the force trend and the }\DIFaddend nearest produced setting \DIFdelbegin \DIFdel{is right beyond }\DIFdelend \DIFaddbegin \DIFadd{decide at 8.2 and }\DIFaddend 8.5 floors\DIFdelbegin \DIFdel{(5.5--10.2); }\DIFdelend \DIFaddbegin \DIFadd{, }\DIFaddend the interval rules \DIFdelbegin \DIFdel{decide }\DIFdelend \DIFaddbegin \DIFadd{(M2, M2n, M3, M3n, M5, M5n) }\DIFaddend only at 13--31\DIFdelbegin \DIFdel{floors, }\DIFdelend \DIFaddbegin \DIFadd{, though they }\DIFaddend are safe at 5.9--16\DIFdelbegin \DIFdel{and at 10 floors send 18--37\,\% of the designs to trial}\DIFdelend . Interpolating in force is much easier than extrapolating \DIFdelbegin \DIFdel{: }\DIFdelend \DIFaddbegin \DIFadd{(}\DIFaddend the nearest setting \DIFdelbegin \DIFdel{decides at }\DIFdelend 4.7 against 9.5 floors\DIFdelbegin \DIFdel{, and M2 and M5 are safe at 2.7 and 2.8 against 9.9 and 8.9 floors. The scaled simulationdecides at 19 floors and}\DIFdelend \DIFaddbegin \DIFadd{), and the scaled simulation}\DIFaddend , extrapolating a fitted trend, falsely accepts failing alternatives as far as 164 floors from their limit.

\emph{A new design family.} \DIFdelbegin \DIFdel{In the two transfers }\DIFdelend \DIFaddbegin \DIFadd{Scored in the floor of the produced family, as the procedure prescribes, }\DIFaddend every rule without the simulation fails \DIFdelbegin \DIFdel{(131--133 }\DIFdelend \DIFaddbegin \DIFadd{in the two transfers (130--135 }\DIFaddend floors), because nothing in one family's production record describes the other. With the simulation, the bias correction decides at \DIFdelbegin \DIFdel{35 floors , }\DIFdelend \DIFaddbegin \DIFadd{33 floors and }\DIFaddend M5 at \DIFdelbegin \DIFdel{38 }\DIFdelend \DIFaddbegin \DIFadd{35, }\DIFaddend and M3 \DIFaddbegin \DIFadd{is safe at 17 floors, M2 }\DIFaddend at \DIFdelbegin \DIFdel{49, and the envelope rule is the safest at 16 floors; no rule is safe closer than that. Scored }\DIFdelend \DIFaddbegin \DIFadd{26 and M5 at 31; }\DIFaddend in the floor of the \DIFdelbegin \DIFdel{produced }\DIFdelend \DIFaddbegin \DIFadd{new }\DIFaddend family, which a \DIFdelbegin \DIFdel{design team would have, M5 decides at 35 and }\DIFdelend \DIFaddbegin \DIFadd{team would not have, M2 }\DIFaddend is safe at \DIFdelbegin \DIFdel{31 floors and the envelope rule }\DIFdelend \DIFaddbegin \DIFadd{16, M3 }\DIFaddend at \DIFdelbegin \DIFdel{41 and 26. Sixteen floors are about }\DIFdelend \DIFaddbegin \DIFadd{21 and M5 at 28. The best rule was thus safe at about 16--17 floors in either floor, some }\DIFaddend 0.8\,mm of mid-side draw-in or 1.3$^\circ$ of wall angle\DIFaddbegin \DIFadd{, but two transfers do not resolve which rule is safest: the safe distance of M3 minus that of M2 lies between $-16$ and $+7$ floors in the produced family's floor and between $-3$ and $+8$ in the new family's}\DIFaddend .

%% BEGIN GENERATED distances
\begin{table}[t]
\tabsize
\DIFdelbeginFL %DIFDELCMD < \setlength{\tabcolsep}{2.2pt}
%DIFDELCMD < %%%
\DIFdelendFL \DIFaddbeginFL \setlength{\tabcolsep}{1.9pt}
\DIFaddendFL \caption{Decisive and safe distances (floors) by the relation of the new design to the produced evidence}\label{tab:distances}
\DIFdelbeginFL %DIFDELCMD < \begin{tabular}{@{}lrrrrrrrrrr@{}}
%DIFDELCMD < \toprule
%DIFDELCMD <  & \multicolumn{2}{c}{New variant} & \multicolumn{4}{c}{New process setting} & \multicolumn{2}{c}{New family} & \multicolumn{2}{c}{Pooled} \\
%DIFDELCMD < \cmidrule(lr){2-3}\cmidrule(lr){4-7}\cmidrule(lr){8-9}\cmidrule(lr){10-11}
%DIFDELCMD < Rule & $\delta_\mathrm{d}/\delta_\mathrm{s}$ & FA & $\delta_\mathrm{d}/\delta_\mathrm{s}$ & FA & interpolated & extrapolated & $\delta_\mathrm{d}/\delta_\mathrm{s}$ & FA & variant & setting \\
%DIFDELCMD < \midrule
%DIFDELCMD < M0 & 108 & 17 & 108 & 17 & 109 & 106 & 108 & 17 & 108 & 108 \\
%DIFDELCMD < M1 & 15 & 15 & 24 & 24 & 11 & 25 & 35 & 33 & 27 & 28 \\
%DIFDELCMD < M1s & 4.9 & 4.7 & 19 & 164 & 19 & 19 & 122 & 122 & 15 & 20 \\
%DIFDELCMD < \quad interval & 4.7--6.3 &  & 14--21 &  & 19--20 & 15--19 & 121--123 &  &  & \\
%DIFDELCMD < M2 & 23/0.15 & 0.30 & 18/8.1 & 8.3 & 18/2.7 & 18/9.9 & 51/16 & 11 & 53/0 & 53/1.0 \\
%DIFDELCMD < M5 & 7.2/0.75 & 0.15 & 19/6.8 & 7.6 & 19/2.8 & 19/8.9 & 38/28 & 32 & 11/0.05 & 28/5.4 \\
%DIFDELCMD < \quad interval & 6.6--8.4 &  & 18--20 &  & 18--20 & 18--20 & 35--39 &  &  & \\
%DIFDELCMD < M5n & 7.2/0.55 & 0.10 & 17/6.1 & 9.0 & 16/0 & 18/8.6 & 132/130 & 167 & 9.8/0.20 & 16/2.7 \\
%DIFDELCMD < \quad interval & 6.4--7.5 &  & 16--18 &  & 15--16 & 16--19 & 130--133 &  &  & \\
%DIFDELCMD < NN & 3.4 & 2.8 & 8.5 & 10 & 4.7 & 9.5 & 131 & 169 & 3.4 & 8.5 \\
%DIFDELCMD < \quad interval & 2.0--4.5 &  & 5.5--10 &  & 4.3--5.1 & 6.6--11 & 129--131 &  &  & \\
%DIFDELCMD < \midrule M0w & 119/94 & 9.2 & 119/94 & 9.2 & 119/94 & 119/77 & 119/94 & 9.2 & 119/94 & 119/94 \\
%DIFDELCMD < Mc & 27 & 28 & 39 & 39 & 14 & 40 & 38 & 42 & 32 & 33 \\
%DIFDELCMD < M3 & 23/0 & 0 & 31/16 & 13 & 29/13 & 33/17 & 49/21 & 20 & 25/0 & 32/5.1 \\
%DIFDELCMD < M1n & 9.1 & 10 & 10 & 12 & 4.8 & 12 & 131 & 173 & 130 & 130 \\
%DIFDELCMD < M2n & 19/0 & 0.20 & 15/5.9 & 9.2 & 13/0 & 16/9.1 & 132/129 & 165 & 325/0 & 325/4.0 \\
%DIFDELCMD < M3n & 11/0.05 & 0 & 13/7.4 & 9.2 & 15/7.0 & 13/7.8 & 133/128 & 164 & 15/0 & 15/4.9 \\
%DIFDELCMD < M6 & 17/0.05 & 0 & 56/6.2 & 5.8 & 47/0 & 56/8.1 & 41/25 & 29 & -- & -- \\
%DIFDELCMD < \botrule
%DIFDELCMD < \end{tabular}
%DIFDELCMD < %%%
\DIFdelendFL \DIFaddbeginFL \begin{tabular}{@{}lrrrrrrrrrr@{}}
\toprule
 & \multicolumn{2}{c}{New variant} & \multicolumn{4}{c}{New process setting} & \multicolumn{2}{c}{New family} & \multicolumn{2}{c}{Pooled} \\
\cmidrule(lr){2-3}\cmidrule(lr){4-7}\cmidrule(lr){8-9}\cmidrule(lr){10-11}
Rule & $\delta_\mathrm{d}/\delta_\mathrm{s}$ & FA & $\delta_\mathrm{d}/\delta_\mathrm{s}$ & FA & interp. & extrap. & $\delta_\mathrm{d}/\delta_\mathrm{s}$ & FA & variant & setting \\
\midrule
M0 & 108 & 17 & 108 & 17 & 109 & 106 & 104 & 28 & 108 & 108 \\
M1 & 15 & 15 & 24 & 24 & 11 & 25 & 33 & 35 & 27 & 28 \\
M1s & 4.9 & 4.7 & 19 & 164 & 19 & 19 & 121 & 125 & 15 & 20 \\
\quad interval & 4.7--6.3 &  & 14--21 &  & 19--20 & 15--19 & 119--121 &  &  & \\
M2 & 23/0.15 & 0.30 & 18/8.1 & 8.3 & 18/2.7 & 18/9.9 & 41/26 & 27 & 53/0 & 53/1.0 \\
\quad interval & 22--23 &  & 17--18 &  & 15--18 & 17--19 & 39--42/9.8--32 &  &  & \\
M5 & 7.2/0.75 & 0.15 & 19/6.8 & 7.6 & 19/2.8 & 19/8.9 & 35/31 & 33 & 11/0.05 & 28/5.4 \\
\quad interval & 6.6--8.4 &  & 18--20 &  & 18--20 & 18--20 & 31--39/28--32 &  &  & \\
M5n & 7.2/0.55 & 0.10 & 17/6.1 & 9.0 & 16/0 & 18/8.6 & 133/128 & 127 & 9.8/0.20 & 16/2.7 \\
\quad interval & 6.4--7.5 &  & 16--18 &  & 15--16 & 16--19 & 131--133/127--128 &  &  & \\
NN & 3.4 & 2.8 & 8.5 & 10 & 4.7 & 9.5 & 130 & 129 & 3.4 & 8.5 \\
\quad interval & 2.0--4.5 &  & 5.5--10 &  & 4.3--5.1 & 6.6--11 & 129--131 &  &  & \\
\midrule M0w & 119/94 & 9.2 & 119/94 & 9.2 & 119/94 & 119/77 & 117/79 & 10 & 119/94 & 119/94 \\
Mc & 27 & 28 & 39 & 39 & 14 & 40 & 43 & 53 & 32 & 33 \\
M3 & 23/0 & 0 & 31/16 & 13 & 29/13 & 33/17 & 40/17 & 18 & 25/0 & 32/5.1 \\
\quad interval & 17--26 &  & 25--36 &  & 24--31 & 25--40 & 36--45/15--19 &  &  & \\
M3n & 11/0.05 & 0 & 13/7.4 & 9.2 & 15/7.0 & 13/7.8 & 135/126 & 124 & 15/0 & 15/4.9 \\
M1n & 9.1 & 10 & 10 & 12 & 4.8 & 12 & 131 & 129 & 130 & 130 \\
M1sn & 3.5 & 3.5 & 8.2 & 9.0 & 4.7 & 8.9 & 131 & 129 & 77 & 111 \\
M2n & 19/0 & 0.20 & 15/5.9 & 9.2 & 13/0 & 16/9.1 & 133/128 & 126 & 325/0 & 325/4.0 \\
M6 & 17/0.05 & 0 & 56/6.2 & 5.8 & 47/0 & 56/8.1 & 35/28 & 28 & -- & -- \\
\botrule
\end{tabular}
\DIFaddendFL \footnotetext{Over all seven characteristics; one number for rules that always decide. FA: safe distance on the false-accept side (failing alternatives). Interval: bootstrap 95\,\% interval of the decisive distance \DIFaddbeginFL \DIFaddFL{and, for a new family, of the safe distance }\DIFaddendFL (fits fixed). 126 cases per relation (\DIFaddbeginFL \DIFaddFL{interp.: }\DIFaddendFL 42 interpolated\DIFdelbeginFL \DIFdelFL{, }\DIFdelendFL \DIFaddbeginFL \DIFaddFL{; extrap.: }\DIFaddendFL 84 extrapolated); a new family rests on two transfers \DIFaddbeginFL \DIFaddFL{and is scored in the floor of the produced family (Section~\ref{sec:floor})}\DIFaddendFL . Pooled: the other geometry's alternatives added to the calibration set. Below the line, the further rules \DIFdelbeginFL \DIFdelFL{; M6}\DIFdelendFL \DIFaddbeginFL \DIFaddFL{(M1sn: the force trend}\DIFaddendFL , the \DIFaddbeginFL \DIFaddFL{counterpart of M1s without the simulation) and, as a robustness variant, M6 (the }\DIFaddendFL M5 mean with a conformal margin on standardised leave-one-out residuals\DIFdelbeginFL \DIFdelFL{, appears only in the robustness analysis}\DIFdelendFL \DIFaddbeginFL \DIFaddFL{)}\DIFaddendFL .}
\end{table}
%% END GENERATED distances

\begin{figure}[t]
\centering
\includegraphics[width=\textwidth]{F3_decisions_compare}
\DIFdelbeginFL %DIFDELCMD < \caption{Verdict shares against the requirement distance $d$ (symmetric logarithmic axis, linear within $\pm$10 floors; 126 cases per panel) with the decisive and safe distances marked}%%%
\DIFdelendFL \DIFaddbeginFL \caption{Verdict shares against the requirement distance $d$ (symmetric logarithmic axis, linear within $\pm$10 floors; 126 cases per panel; a new family in the produced family's floor), with the decisive and safe distances}\DIFaddendFL \label{fig:decisions}
\end{figure}

\subsection{Where the distances come from}\label{sec:res-drivers}

Three features of the \DIFdelbegin \DIFdel{evaluation design shape the }\DIFdelend \DIFaddbegin \DIFadd{design shape these }\DIFaddend distances. First, the sibling relation mixes replication with variation: in 58 of the 126 \DIFdelbegin \DIFdel{new-variant }\DIFdelend cases neither sibling differs resolvably from the held-out alternative, and the nearest \DIFdelbegin \DIFdel{produced }\DIFdelend setting is then right beyond 0.85 floors, the reproducibility of $q_{95}$ between near-replicate runs; \DIFdelbegin \DIFdel{when both differ }\DIFdelend \DIFaddbegin \DIFadd{it needs 2.5 floors when one sibling differs }\DIFaddend resolvably (34 cases) \DIFdelbegin \DIFdel{, it needs }\DIFdelend \DIFaddbegin \DIFadd{and }\DIFaddend 4.9 \DIFdelbegin \DIFdel{floors }\DIFdelend \DIFaddbegin \DIFadd{when both do (34 cases), where the force trend needs 4.5 }\DIFaddend and M5 9.7. Second, the extrapolation \DIFdelbegin \DIFdel{cases are }\DIFdelend \DIFaddbegin \DIFadd{is }\DIFaddend confounded with stroke speed: the \DIFdelbegin \DIFdel{production-only rules lose almost only }\DIFdelend \DIFaddbegin \DIFadd{nearest setting decides at 11 floors }\DIFaddend at 100\,kN, where the press also ran slower\DIFdelbegin \DIFdel{(the nearest setting decides at 11 floors }\DIFdelend \DIFaddbegin \DIFadd{, }\DIFaddend against 4.9 at 500\,kN \DIFdelbegin \DIFdel{), whereas the simulation-based interval rules are unsafe at bothends. Third, the production record of }\DIFdelend \DIFaddbegin \DIFadd{(the force trend at 8.9 at both). Taken together, the nearest setting decided at about 5 floors whenever the new design differed resolvably and force and speed did not change together, though each stratum has 34 to 58 cases and each distance rests on a handful of verdicts. Third, }\DIFaddend the other family\DIFaddbegin \DIFadd{'s record }\DIFaddend does not help \DIFdelbegin \DIFdel{within a family}\DIFdelend \DIFaddbegin \DIFadd{with a produced sibling}\DIFaddend : pooling it makes most rules less decisive (M5 from 7.2 to 11 floors\DIFdelbegin \DIFdel{for a new variant, M2 from 23 to 53}\DIFdelend \DIFaddbegin \DIFadd{)}\DIFaddend , \DIFdelbegin \DIFdel{M2n from 19 to 325), although it raises the coverage of }\DIFdelend \DIFaddbegin \DIFadd{though it raises }\DIFaddend M5\DIFaddbegin \DIFadd{'s coverage }\DIFaddend from 86 to 90\,\%\DIFdelbegin \DIFdel{and makes the centre of its interval the most accurate point prediction within the family (2.9 floors against 3.4 for the nearest setting)}\DIFdelend .

The relation therefore does not govern decision fitness in general. Over the eleven calibrated rules and three relations, the relation explains \DIFdelbegin \DIFdel{69}\DIFdelend \DIFaddbegin \DIFadd{67}\DIFaddend \,\% of the variation of the log decisive distance and the rule 6\,\% \DIFdelbegin \DIFdel{, but this share comes from }\DIFdelend \DIFaddbegin \DIFadd{(69 and 6\,\% in }\DIFaddend the new family\DIFdelbegin \DIFdel{, where the }\DIFdelend \DIFaddbegin \DIFadd{'s own floor), but only because }\DIFaddend rules without the simulation fail \DIFaddbegin \DIFadd{for a new family }\DIFaddend by construction; \DIFaddbegin \DIFadd{over the two relations }\DIFaddend within a family\DIFaddbegin \DIFadd{, }\DIFaddend the rule explains 68\,\% and the relation 15\,\% \DIFdelbegin \DIFdel{. The }\DIFdelend \DIFaddbegin \DIFadd{(71 and 15\,\% with M1sn). Pooled over strata, the }\DIFaddend best attainable distance \DIFdelbegin \DIFdel{still grows with the noveltyof the design, from }\DIFdelend \DIFaddbegin \DIFadd{grows with novelty: }\DIFaddend 3.4 floors with a \DIFdelbegin \DIFdel{sibling to 8.5 }\DIFdelend \DIFaddbegin \DIFadd{produced sibling, 8.2 }\DIFaddend for a new setting\DIFdelbegin \DIFdel{and 35 }\DIFdelend \DIFaddbegin \DIFadd{, 33 }\DIFaddend for a new family.

\subsection{What the simulation and learning contribute}\label{sec:res-ml}

Pairing each calibrated rule with its counterpart without the simulation isolates what the simulation adds as \DIFdelbegin \DIFdel{it is }\DIFdelend used. Added as an offset, it adds nothing \DIFdelbegin \DIFdel{within a family }\DIFdelend \DIFaddbegin \DIFadd{with a produced sibling }\DIFaddend and degrades two rules: M5 decides where M5n does (difference 0.0 floors, $-0.4$ to $+1.1$), \DIFdelbegin \DIFdel{and }\DIFdelend \DIFaddbegin \DIFadd{while }\DIFaddend M2 and M1 decide 3.9 ($+0.5$ to $+7.4$) and 6.2 floors ($+4.6$ to $+9.7$) further out than M2n and M1n, \DIFaddbegin \DIFadd{and M1 13.5 further out for a new setting, }\DIFaddend because the simulated force effect is too large. Rescaled, it \DIFdelbegin \DIFdel{adds information within a family, where the scaled simulation decides 4.2 }\DIFdelend \DIFaddbegin \DIFadd{does no better: within a geometry the nominal simulation is a function of the force, and the straight line in force (M1sn) decides 1.4 }\DIFaddend floors closer than the \DIFdelbegin \DIFdel{median of the produced alternatives ($-5.4$ to $-0.5$) , though 1.5 floors less close than the nearest setting; }\DIFdelend \DIFaddbegin \DIFadd{scaled simulation with a sibling ($+0.65$ to $+3.0$) and 11 closer }\DIFaddend for a new setting \DIFdelbegin \DIFdel{the rescaled trend is worse than none}\DIFdelend \DIFaddbegin \DIFadd{($+5.3$ to $+13.8$), in every resample and refit configuration; the 4.2 floors that the scaled simulation gains over the median of the produced alternatives with a sibling come from the force level}\DIFaddend . Only for a new family is the simulation decisive, by \DIFdelbegin \DIFdel{about 94 floors }\DIFdelend \DIFaddbegin \DIFadd{92 to 98 floors as an offset and 10 rescaled}\DIFaddend , and even there production knowledge carries over where a drawing nominal transfers: for the wall angle, the design angle plus the produced family's deviation from its own \DIFdelbegin \DIFdel{design angle decides at 2.0 }\DIFdelend \DIFaddbegin \DIFadd{decides at 1.9 }\DIFaddend floors, against \DIFdelbegin \DIFdel{7.3 }\DIFdelend \DIFaddbegin \DIFadd{7.5 }\DIFaddend for M5 \DIFdelbegin \DIFdel{. A }\DIFdelend \DIFaddbegin \DIFadd{(for the arm angle and the dome, whose nominal is zero in both families, this baseline is M1n, at 44 and 27 floors). A Gaussian-process }\DIFaddend surrogate of the simulations \DIFdelbegin \DIFdel{reproduces them to within 1.8 }\DIFdelend \DIFaddbegin \DIFadd{over force, friction and thickness, with a leave-one-out root-mean-square error below 2 }\DIFaddend floors for 12 of the 14 characteristics and geometries, \DIFdelbegin \DIFdel{yet its decision fitness is that of the simulation it reproduces}\DIFdelend \DIFaddbegin \DIFadd{decides as the simulation does (109 against 108 floors as M0, 15 against 15 as M1 with a sibling)}\DIFaddend .

Learning adds intervals rather than accuracy. Used as point rules, the centres of the M5, M5n and M3n intervals decide at 3.3, 3.1 and 3.7 floors \DIFdelbegin \DIFdel{within the family}\DIFdelend \DIFaddbegin \DIFadd{with a produced sibling}\DIFaddend , as close as the nearest setting \DIFdelbegin \DIFdel{, so the ranking of the learned rules rests on the width of their intervals. With the noise floor set by the variation between runs, }\DIFdelend \DIFaddbegin \DIFadd{(Table~\ref{tab:intervals}). With a sibling }\DIFaddend the Gaussian-process intervals cover 85--86\,\% of the true $q_{95}$ \DIFdelbegin \DIFdel{within the family }\DIFdelend against a nominal 90\,\%, \DIFaddbegin \DIFadd{but their noise variance sits at its lower bound, the variance between the calibration series of one setting, in 97--99\,\% of these fits and in every new-setting and new-family fit, so this coverage comes from the quasi-replicate series rather than from the learner; }\DIFaddend the largest-residual margin \DIFaddbegin \DIFadd{covers }\DIFaddend 89\,\% \DIFdelbegin \DIFdel{(it guarantees at least $(n-1)/(n+1)=7/9$ under exchangeability, which these cases lack) }\DIFdelend and the jackknife+ \DIFdelbegin \DIFdel{interval of M3 }\DIFdelend \DIFaddbegin \DIFadd{intervals }\DIFaddend 96\,\% \DIFdelbegin \DIFdel{; for }\DIFdelend \DIFaddbegin \DIFadd{(M3) and 89\,\% (M3n), above their guarantee of 7/9. For }\DIFaddend a new setting the calibrated intervals cover 43--55\,\% \DIFdelbegin \DIFdel{, and }\DIFdelend \DIFaddbegin \DIFadd{(M2n and M5n 93\,\% interpolating), }\DIFaddend for a new family the Gaussian processes at most 12\,\%\DIFdelbegin \DIFdel{. The inductive biases explain much of this}\DIFdelend : trees do not extrapolate, two force levels identify at most a linear trend, and in the median new-setting fit the length scale of the force sits at its lower bound\DIFdelbegin \DIFdel{, so the Gaussian process reverts to its mean}\DIFdelend .

\subsection{Calibration budget}\label{sec:res-budget}

\DIFdelbegin \DIFdel{Figure~\ref{fig:budget} shows how the distances depend on the number of produced alternatives of the family and on their relation to the new design. }\DIFdelend A single produced sibling decides \DIFaddbegin \DIFadd{(Fig.~\ref{fig:budget})}\DIFaddend : with one alternative at the new design's force, the bias correction, the envelope rule and the nearest setting coincide \DIFdelbegin \DIFdel{by construction (the simulation terms cancel and a single residual gives no margin) }\DIFdelend and decide at 3.5 floors (2.3--4.6) without a wrong verdict at 10 floors. Further alternatives do not improve the nearest setting; they make the bias correction worse (12--18 floors), the envelope rule safer but less decisive, and \DIFdelbegin \DIFdel{the Gaussian processes }\DIFdelend \DIFaddbegin \DIFadd{M5, shown from three alternatives because a Gaussian process fits its hyperparameters to them, }\DIFaddend safer and more decisive (\DIFdelbegin \DIFdel{from 12--14 floors with two alternatives }\DIFdelend \DIFaddbegin \DIFadd{13 }\DIFaddend to 7.2 \DIFdelbegin \DIFdel{with eight}\DIFdelend \DIFaddbegin \DIFadd{floors}\DIFaddend ). Without a sibling, bracketing the new force buys safety but not decisiveness: \DIFdelbegin \DIFdel{from two alternatives on, the interval rules }\DIFdelend \DIFaddbegin \DIFadd{the envelope rule and M5 }\DIFaddend are safe at \DIFdelbegin \DIFdel{0--4}\DIFdelend \DIFaddbegin \DIFadd{2.7--4}\DIFaddend .8 floors \DIFdelbegin \DIFdel{with no wrong verdict at 10 floors }\DIFdelend but decide only at 15--25\DIFdelbegin \DIFdel{floors}\DIFdelend , against 4.7--6.0 for the nearest setting\DIFdelbegin \DIFdel{. Extrapolation }\DIFdelend \DIFaddbegin \DIFadd{, and extrapolation }\DIFaddend stays unsafe at every budget. Averaged over all subsets, the distances shrink with the budget mainly because the share of subsets containing a sibling grows, from 46\,\% with two alternatives to 96\,\% with six\DIFdelbegin \DIFdel{; the odd--even pattern of the bias correction comes from its median offset, which averages two forces only for an even number of alternatives}\DIFdelend . With calibration series of 50, 100 or 250 parts instead of 500, no distance within a family moves by more than \DIFdelbegin \DIFdel{1.2 }\DIFdelend \DIFaddbegin \DIFadd{1.5 }\DIFaddend floors, and for a new family by at most 4 \DIFdelbegin \DIFdel{floors }\DIFdelend (Table~\ref{tab:robust}): \DIFdelbegin \DIFdel{the qualification needs short runsof several variants}\DIFdelend \DIFaddbegin \DIFadd{rules can be calibrated on short runs, here scored against the truths and floors of the full series}\DIFaddend , and only the floor needs a run long enough to show the drift.

\begin{figure}[t]
\centering
\includegraphics[width=\textwidth]{F4_budget}
\caption{Decisive and safe distances against the number of produced alternatives of the family, by the relation of the new design to them; bands: 95\,\% intervals over the held-out alternatives}\label{fig:budget}
\end{figure}

\subsection{The procedure's decision rule}\label{sec:res-procedure}

What a designer observes for a single verdict is the margin between the predicted interval and the requirement, and its value as evidence depends on the relation (Fig.~\ref{fig:reliability}): \DIFdelbegin \DIFdel{for a new variant }\DIFdelend \DIFaddbegin \DIFadd{with a produced sibling }\DIFaddend at least 95\,\% of the decided verdicts of the interval rules are correct at every margin\DIFdelbegin \DIFdel{; }\DIFdelend \DIFaddbegin \DIFadd{, }\DIFaddend for a new setting the share grows with the margin, \DIFdelbegin \DIFdel{for the nearest setting from 85\,\% to 95\,\% beyond 5 floors; }\DIFdelend \DIFaddbegin \DIFadd{and }\DIFaddend for a new family \DIFaddbegin \DIFadd{M2 and }\DIFaddend M5 \DIFdelbegin \DIFdel{is right in 71--74}\DIFdelend \DIFaddbegin \DIFadd{are right in at most 79}\DIFaddend \,\% of \DIFdelbegin \DIFdel{its decided verdictswhatever the margin.
}\DIFdelend \DIFaddbegin \DIFadd{their decided verdicts.
}

\DIFaddend Step 5 turns this into guard bands \DIFdelbegin \DIFdel{: none for }\DIFdelend \DIFaddbegin \DIFadd{(Table~\ref{tab:step6}). Under the reference prior }\DIFaddend the interval rules \DIFdelbegin \DIFdel{within the family and 0.25 floors for }\DIFdelend \DIFaddbegin \DIFadd{need none with a sibling and }\DIFaddend the nearest setting \DIFaddbegin \DIFadd{0.25 floors}\DIFaddend ; \DIFdelbegin \DIFdel{2.0 floors}\DIFdelend for an interpolated setting \DIFdelbegin \DIFdel{, with which }\DIFdelend the nearest setting \DIFaddbegin \DIFadd{needs 2.0 floors and then }\DIFaddend decides at 6.7 floors \DIFdelbegin \DIFdel{, is safe at 2.7 and sends no design to trialat }\DIFdelend \DIFaddbegin \DIFadd{without trials at }\DIFaddend 10 \DIFdelbegin \DIFdel{floors; 7.0 floors }\DIFdelend \DIFaddbegin \DIFadd{floors, whereas M2 and M5 need 0.75 and 0.5 floors but send 31 and 40\,\% to trial; }\DIFaddend for an extrapolated \DIFdelbegin \DIFdel{one, with which it decides at 17 }\DIFdelend \DIFaddbegin \DIFadd{setting it needs 7.0 }\DIFaddend floors and sends a quarter \DIFdelbegin \DIFdel{of the designs to trialat 10 floors; and }\DIFdelend \DIFaddbegin \DIFadd{to trial. M0 and M0w never qualify, nor do M1 and Mc with a sibling, M1s for an interpolated setting, M1, M1s, M2, M5 and Mc for an extrapolated one, or any rule }\DIFaddend for a new family\DIFdelbegin \DIFdel{no guard band up to 20 floors reaches 95\,\%, so every }\DIFdelend \DIFaddbegin \DIFadd{.
}

\DIFadd{These bands belong to the prior as much as to the rules. A prior concentrated near the truth, where release-level requirements lie (the grid within $\pm5$ floors), raises the band of the nearest setting to 2.0 floors with a sibling, 4.25 interpolating and 14.25 extrapolating; a uniform prior within $\pm20$ floors lowers them to 0, 0.25 and 3.5 floors and qualifies M2 and M5 for an extrapolated setting (6.0 and 5.0 floors), and one within $\pm50$ floors qualifies M2, M3 and M5 for a new family (6.25, 6.5 and 15 floors). That every new-family }\DIFaddend design goes to trial \DIFdelbegin \DIFdel{. The guard }\DIFdelend \DIFaddbegin \DIFadd{thus follows from the reference prior and the 20-floor cap, as the practice of trying out every new die would have it. The }\DIFaddend bands are chosen on the evaluation data, which makes the trial rates optimistic\DIFdelbegin \DIFdel{. For a new setting the range guard refuses exactly the extrapolated forces, after which M2, M5 and the nearest setting make no wrong verdict at 10 floors; the novelty guardalso refuses many cases that }\DIFdelend \DIFaddbegin \DIFadd{, and the guards refuse indiscriminately: 96\,\% of the range guard's refusals for a new setting }\DIFaddend would have been \DIFdelbegin \DIFdel{decided correctly}\DIFdelend \DIFaddbegin \DIFadd{right for the nearest setting and 72\,\% for M5 (novelty guard 96 and 73\,\%)}\DIFaddend .

\emph{A worked example.} A tryout engineer has to decide whether the concave cup at 300\,kN with medium lubrication, a new \DIFdelbegin \DIFdel{process }\DIFdelend setting, can be released without a trial run, with the six alternatives at 100 and 500\,kN produced and a requirement\DIFaddbegin \DIFadd{, chosen for illustration, }\DIFaddend that 95\,\% of the parts have a mid-side draw-in\DIFaddbegin \DIFadd{, by which tryout sets the blank-holder force, }\DIFaddend of at most 15.30\,mm. The floor is 0.050\,mm, the relation is interpolation and no guard fires. \DIFdelbegin \DIFdel{The }\DIFdelend \DIFaddbegin \DIFadd{Alone, the }\DIFaddend nominal simulation predicts 15.66\,mm and fails the design. The nearest produced setting\DIFaddbegin \DIFadd{, the mean of the six alternatives at the neighbouring forces, }\DIFaddend predicts 14.99\,mm, a margin of 6.2 floors, more than its guard band of 2.0 floors \DIFaddbegin \DIFadd{(5.0 under the grid within $\pm3$ floors)}\DIFaddend , so step 6 returns meets\DIFdelbegin \DIFdel{; beyond a margin of 6 floors every decided verdict of this rule for an interpolated setting was correct. }\DIFdelend \DIFaddbegin \DIFadd{. Widened by their guard bands, }\DIFaddend M2, M5 and M5n return \DIFdelbegin \DIFdel{intervals reaching 15.37, 15.58 and 15.33}\DIFdelend \DIFaddbegin \DIFadd{14.41--15.41, 14.23--15.61 and 14.66--15.33}\DIFaddend \,mm and send the design to trial. The 500 parts had a 95th percentile of 14.88\,mm, 8.4 floors below the limit\DIFdelbegin \DIFdel{: the procedure was right, the simulation wrong and the interval rules cautious}\DIFdelend \DIFaddbegin \DIFadd{. The case is typical, since at 8.5 floors above the truth step 6 with the nearest setting is right in all 42 interpolated cases, but its guard band was estimated on data that include it}\DIFaddend .

\begin{figure}[t]
\centering
\includegraphics[width=\textwidth]{F5_reliability}
\DIFdelbeginFL %DIFDELCMD < \caption{Share of correct verdicts among the decided ones against the observable margin between interval and requirement, for requirements within 20 floors of the truth}%%%
\DIFdelendFL \DIFaddbeginFL \caption{Share of correct verdicts among the decided ones with at least the given margin between interval and requirement, under the reference prior (47 grid distances within $\pm$20 floors, equal weights); filled markers: guard bands}\DIFaddendFL \label{fig:reliability}
\end{figure}

\subsection{Requirements anchored in capability and in tolerances}\label{sec:res-requirements}

Upper limits at which an alternative has a performance index of 1.0, 1.33, 1.67 or 2.0 lie a median 1.0, 1.7, 2.5 and 3.1 floors above its 95th percentile\DIFdelbegin \DIFdel{, }\DIFdelend \DIFaddbegin \DIFadd{; up to an index of 1.33 they lie }\DIFaddend inside every decisive distance \DIFdelbegin \DIFdel{; }\DIFdelend \DIFaddbegin \DIFadd{except a near-replicate sibling's (0.85 floors). The test is one-sided by construction: }\DIFaddend every alternative meets \DIFdelbegin \DIFdel{them}\DIFdelend \DIFaddbegin \DIFadd{these limits}\DIFaddend , so a verdict \DIFdelbegin \DIFdel{is }\DIFdelend \DIFaddbegin \DIFadd{can only be }\DIFaddend right, a false reject or a trial\DIFaddbegin \DIFadd{, and a rule biased towards meets scores well}\DIFaddend . At an index of 1.33 the \DIFdelbegin \DIFdel{nearest setting still says meets in 91\,\% of the new-variant cases, because most of its errors arefar smaller than its 95\,\% distance, whereas M5 says so in 44\,\% and sends the rest to trial; }\DIFdelend \DIFaddbegin \DIFadd{shares of right verdicts, false rejects and trials are, with a produced sibling and }\DIFaddend for a new setting\DIFdelbegin \DIFdel{the two are right in }\DIFdelend \DIFaddbegin \DIFadd{, 91/9/0 and }\DIFaddend 65\DIFaddbegin \DIFadd{/35/0\,\% for the nearest setting, 44/2/53 }\DIFaddend and 32\DIFaddbegin \DIFadd{/16/52}\DIFaddend \,\% \DIFdelbegin \DIFdel{, }\DIFdelend \DIFaddbegin \DIFadd{for M5, and 81/19/0 }\DIFaddend and \DIFaddbegin \DIFadd{79/21/0\,\% for the scaled simulation, which falsely accepts failing alternatives elsewhere; }\DIFaddend for a new family no rule \DIFdelbegin \DIFdel{exceeds }\DIFdelend \DIFaddbegin \DIFadd{is right in more than }\DIFaddend 54\,\%. Release-level requirements are \DIFdelbegin \DIFdel{therefore }\DIFdelend \DIFaddbegin \DIFadd{thus }\DIFaddend decided at the design stage only \DIFdelbegin \DIFdel{within a family}\DIFdelend \DIFaddbegin \DIFadd{with a produced sibling}\DIFaddend , and then by the production record. General angular tolerances of ISO~2768-1 \citep{iso1989tolerances} give a coarser scenario\DIFdelbegin \DIFdel{that mostly tests whether a rule recognises gross failures (}\DIFdelend \DIFaddbegin \DIFadd{: only }\DIFaddend 32 of its 108 decisions lie within 5 floors of the limit\DIFdelbegin \DIFdel{)}\DIFdelend \DIFaddbegin \DIFadd{, 6 of them on failing alternatives, and a cost comparison inherits this coarseness}\DIFaddend . With a false reject at half the cost of a false accept, \DIFdelbegin \DIFdel{its }\DIFdelend \DIFaddbegin \DIFadd{the }\DIFaddend cheapest rule is \DIFaddbegin \DIFadd{(i) with a produced sibling, }\DIFaddend the nearest setting\DIFdelbegin \DIFdel{at any costfor a new variant; }\DIFdelend \DIFaddbegin \DIFadd{, the scaled simulation and the force trend, tied at zero cost; (ii) }\DIFaddend for a new setting\DIFaddbegin \DIFadd{, }\DIFaddend the \DIFdelbegin \DIFdel{production-only interval rules while a trial costs at most a tenth of a false accept, and the nearest setting beyond; }\DIFdelend \DIFaddbegin \DIFadd{force trend at every trial cost; and (iii) }\DIFaddend for a new family\DIFaddbegin \DIFadd{, }\DIFaddend a trial of every design while a trial costs at most a tenth \DIFaddbegin \DIFadd{of a false accept}\DIFaddend , the envelope rule at a fifth and M5 beyond.

\subsection{Robustness}\label{sec:res-robust}

\DIFdelbegin \DIFdel{The analysis was repeated }\DIFdelend \DIFaddbegin \DIFadd{Table~\ref{tab:robust} repeats the analysis }\DIFaddend under variants of the data, \DIFdelbegin \DIFdel{the constants and the model specifications(Table~\ref{tab:robust})}\DIFdelend \DIFaddbegin \DIFadd{constants and model specifications}\DIFaddend . Adjusting for punch temperature changes no calibrated distance by more than 2.6 floors. Dropping the warm-up of the first 150 parts shrinks the floors and so lengthens the distances\DIFdelbegin \DIFdel{(the nearest setting 4.1 floors with a sibling, 12 for a new setting)}\DIFdelend , and floors taken between series \DIFdelbegin \DIFdel{, 1.5 to 9.1 times the reference, shorten them(0.75 and 2.7 floors)}\DIFdelend \DIFaddbegin \DIFadd{shorten them}\DIFaddend . Other floor protocols rescale the \DIFaddbegin \DIFadd{decisive }\DIFaddend distances by 0.60 to 1.42, \DIFaddbegin \DIFadd{and }\DIFaddend the Gaussian-process specifications move them by at most \DIFdelbegin \DIFdel{2.1 floors , and in the unit of the reproducibility of $q_{95}$ the nearest setting decides at 2.8 }\DIFdelend \DIFaddbegin \DIFadd{1.2 floors }\DIFaddend and \DIFdelbegin \DIFdel{6.9 floors}\DIFdelend \DIFaddbegin \DIFadd{the safe distances by at most 2.1}\DIFaddend . In every variant that includes \DIFdelbegin \DIFdel{it, the nearest setting remains the }\DIFdelend \DIFaddbegin \DIFadd{them, the }\DIFaddend most decisive rule within a family \DIFdelbegin \DIFdel{and the envelope rule the safest for a new family}\DIFdelend \DIFaddbegin \DIFadd{is the nearest setting or the force trend, and with a produced sibling the nearest setting (tied once)}\DIFaddend . The refitting bootstrap widens the intervals of the learned rules most (M5 4.0--11 floors \DIFdelbegin \DIFdel{for a new variant}\DIFdelend \DIFaddbegin \DIFadd{with a sibling}\DIFaddend , the nearest setting 2.3--4.6), yet \DIFaddbegin \DIFadd{in all 200 resamples (48 configurations) }\DIFaddend the nearest setting is \DIFdelbegin \DIFdel{the more decisive in all 200 resamples, against }\DIFdelend \DIFaddbegin \DIFadd{more decisive than }\DIFaddend M5 and M5n\DIFaddbegin \DIFadd{, with a sibling }\DIFaddend and for a new setting alike\DIFaddbegin \DIFadd{, and the force trend than the scaled simulation (by 1.35--3.6 and 8.8--12 floors); point estimates at the edge of their refit intervals (M1s with a sibling at 4.9 floors, 4.9--7.3, and the safe distances of M2 and M5) reflect what losing a pattern does}\DIFaddend .

%% ----------------------------------------------------------------------
\section{Discussion}\label{sec:discussion}

\subsection{What the evaluation establishes}

Five findings answer the research questions. Two follow from the definitions and hold wherever the framework is applied: the floor puts every rule and characteristic on a common production-referenced scale, though not on a common scale of nonconforming fractions, and as a within-run quantity it \DIFdelbegin \DIFdel{is a lower bound of }\DIFdelend \DIFaddbegin \DIFadd{bounds }\DIFaddend the variation of series production \DIFaddbegin \DIFadd{from below }\DIFaddend (RQ1); and an interval rule trades decisiveness for safety, which the two distances make explicit, \DIFdelbegin \DIFdel{and reaches its nominal }\DIFdelend \DIFaddbegin \DIFadd{while its distribution-free intervals are guaranteed their }\DIFaddend level only where the calibration alternatives are exchangeable with the new design. Three are observations on this dataset and hypotheses for confirmation: the best attainable distance grows with the novelty of the design, \DIFaddbegin \DIFadd{from 3.4 floors with a produced sibling (the nearest setting 0.85 for near-replicates, 4.9 for resolvable siblings) to 8.2 for a new setting (4.7 interpolated; 4.9 and 8.9 extrapolated to 500 and 100\,kN) and 33 for a new family, }\DIFaddend while within a family the choice of rule matters more than the relation (RQ2); production evidence buys safety before decisiveness, \DIFdelbegin \DIFdel{since interval rules become safe with a few produced alternatives that bracket the new design and decisive only with near-replicates}\DIFdelend \DIFaddbegin \DIFadd{and interval rules held their level with a produced sibling but not when extrapolating or for a new family}\DIFaddend ; and with this simulator and six to nine produced \DIFdelbegin \DIFdel{settings per fit}\DIFdelend \DIFaddbegin \DIFadd{alternatives at two or three force levels}\DIFaddend , the simulation\DIFdelbegin \DIFdel{contributed }\DIFdelend \DIFaddbegin \DIFadd{, }\DIFaddend as an offset \DIFaddbegin \DIFadd{or rescaled, contributed }\DIFaddend only across families, \DIFdelbegin \DIFdel{a fitted scale made it useful within a family but harmful for a new setting, }\DIFdelend and learned models predicted as well as the nearest produced setting \DIFaddbegin \DIFadd{with a produced sibling }\DIFaddend while adding abstention (RQ3).

\begin{table}[t]
\tabsize
\setlength{\tabcolsep}{4pt}
\caption{Guidance for design-stage manufacturability decisions, with its basis}\label{tab:guidelines}
\DIFdelbeginFL %DIFDELCMD < \begin{tabular}{@{}L{2.9cm}L{7.75cm}L{1.85cm}@{}}
%DIFDELCMD < \toprule
%DIFDELCMD < Situation & Guidance & Basis \\
%DIFDELCMD < \midrule
%DIFDELCMD < \multicolumn{3}{@{}l}{\emph{From the framework (general)}} \\
%DIFDELCMD < Before trusting a predictor & Measure the floor on consecutive production of the family with a stated protocol and report $\sigma_\mathrm{b}$, $\sigma_\mathrm{w}$ and its lag dependence; treat it as a lower bound of series variation. & definitions \\
%DIFDELCMD < Qualifying a rule & Evaluate it for the evidence relation of the coming designs by grouped cross-validation, and split near-replicate siblings from resolvable ones. & definitions \\
%DIFDELCMD < Choosing a rule & Compare it with the nearest produced setting and with its own version without the simulation, on decisive and safe distance jointly or by expected cost. & definitions \\
%DIFDELCMD < Deciding a single design & Apply the rule's guard band to the observable margin; a margin inside the guard band, or a fired guard, means a trial. & definitions \\
%DIFDELCMD < Fixing a requirement & If the expected margin to the capability is below the safe distance of the best qualified rule, plan a trial; no decision is reliable closer than a trial reproduces the quantile. & definitions \\
%DIFDELCMD < \multicolumn{3}{@{}l}{\emph{Observed on this dataset (18 alternatives, 2 geometries, 7 characteristics)}} \\
%DIFDELCMD < New variant with a produced sibling & nearest produced setting right beyond 3.4 floors (0.85 for a near-replicate); calibrated interval rules safe below one floor & 126 cases \\
%DIFDELCMD < New process setting & nearest produced setting right beyond 4.7 floors interpolated and 9.5 extrapolated; interval rules safe at 5.9--16 floors & 42 + 84 cases \\
%DIFDELCMD < New family & best rule safe at 16 floors; rules without the simulation fail; a transferable drawing nominal helps & two transfers \\
%DIFDELCMD < Release-level requirements & lie within 1--3 floors of the 95th percentile, inside every decisive distance & 126 cases \\
%DIFDELCMD < \botrule
%DIFDELCMD < \end{tabular}
%DIFDELCMD < %%%
\DIFdelendFL \DIFaddbeginFL \begin{tabular}{@{}L{2.9cm}L{7.75cm}L{1.85cm}@{}}
\toprule
Situation & Guidance & Basis \\
\midrule
\multicolumn{3}{@{}l}{\emph{From the framework (prescriptive; usefulness not evaluated)}} \\
Before relying on a predictor & Measure the floor on consecutive production of the family with a stated protocol and report $\sigma_\mathrm{b}$, $\sigma_\mathrm{w}$ and its lag dependence; treat it as a lower bound of series variation. & framework \\
Qualifying a rule & Evaluate it for the evidence relation of the coming designs by grouped cross-validation, and split near-replicate siblings from resolvable ones. & framework \\
Choosing a rule & Compare it with the nearest produced setting and with its own version without the simulation, on decisive and safe distance jointly or by expected cost. & framework \\
Deciding a single design & Derive the rule's guard band for the distribution of requirement distances the team expects and apply it to the observable margin; a margin inside it, or a fired guard, means a trial. & framework \\
Planning a requirement & As a heuristic, plan a trial if the expected margin between requirement and predicted 95th percentile is below the safe distance of the best qualified rule; no decision is reliable closer than a trial reproduces the quantile. & framework \\
\multicolumn{3}{@{}l}{\emph{Observed on this dataset (18 alternatives, 2 geometries, 7 characteristics)}} \\
New variant with a produced sibling & nearest produced setting right beyond 3.4 floors (0.85 for near-replicates, 4.9 for resolvable siblings); interval rules safe below one floor, coverage 85--96\,\% & 126 cases \\
New process setting & nearest setting or force trend right beyond 4.7 floors interpolated, 4.9 extrapolated to 500\,kN and 8.9 to 100\,kN (with a speed change); interval rules safe at 5.9--16 floors, coverage 43--55\,\% & 42 + 84 cases \\
New family & best rule safe at 16--17 floors in either family's floor, which rule not resolved; rules without the simulation fail; a transferable drawing nominal helps & two transfers \\
Release-level requirements & lie 1.7 floors above the 95th percentile at an index of 1.33, inside every decisive distance except a near-replicate's; decided only with a produced sibling & 126 cases \\
\botrule
\end{tabular}
\DIFaddendFL \end{table}

\subsection{Using the framework in design}

\DIFaddbegin \DIFadd{As a planning heuristic (}\DIFaddend Table~\ref{tab:guidelines}\DIFdelbegin \DIFdel{separates guidance that follows from the framework from observations on this dataset. Before a requirement is fixed}\DIFdelend \DIFaddbegin \DIFadd{)}\DIFaddend , the expected margin between \DIFdelbegin \DIFdel{requirement and capability }\DIFdelend \DIFaddbegin \DIFadd{a requirement and the predicted 95th percentile }\DIFaddend can be compared with the safe distance of the best rule qualified for the relation\DIFdelbegin \DIFdel{of the design; if the margin }\DIFdelend \DIFaddbegin \DIFadd{; if it }\DIFaddend is smaller, a trial or \DIFdelbegin \DIFdel{further }\DIFdelend \DIFaddbegin \DIFadd{more }\DIFaddend production evidence has to be planned\DIFaddbegin \DIFadd{, and the decision itself needs the guard band, because the safe distance is defined on the true margin}\DIFaddend . Release-level requirements make this the usual case\DIFdelbegin \DIFdel{: they lie within a few floors of the capability, where only a near-replicate production record decides reliably, consistent with }\DIFdelend \DIFaddbegin \DIFadd{, as }\DIFaddend the practice of \DIFdelbegin \DIFdel{requiring }\DIFdelend a production trial run for a process change \DIFdelbegin \DIFdel{. For a single design, step 6 gives a verdict whose guard band a designer can see, and the worked example shows that the most decisive rule is also the most transparent: the nearest produced setting points to the produced alternative it copies, whereas a Gaussian-process interval does not. Which rule to prefer depends on the costs, which differ by relation: a trial of }\DIFdelend \DIFaddbegin \DIFadd{expects: the rules decide reliably where a trial is cheapest, }\DIFaddend a new force on an existing tool \DIFdelbegin \DIFdel{is }\DIFdelend \DIFaddbegin \DIFadd{being }\DIFaddend a short press run, \DIFdelbegin \DIFdel{a trial of }\DIFdelend \DIFaddbegin \DIFadd{and none does near release-level limits where a trial is dearest, for }\DIFaddend a new family \DIFaddbegin \DIFadd{that }\DIFaddend may need a new tool. \DIFdelbegin \DIFdel{In a company, }\DIFdelend \DIFaddbegin \DIFadd{The most decisive rule is arguably also the most transparent: the nearest produced setting points to the produced alternatives whose record it averages. The variants that }\DIFaddend steps 1 to 5 \DIFdelbegin \DIFdel{fit the start of production or the }\DIFdelend \DIFaddbegin \DIFadd{need come from process development and tryout, whose short series of several settings support the sibling and new-setting relations, and from a family's production history across part numbers, which supports the new-family relation; }\DIFaddend production part approval\DIFdelbegin \DIFdel{of a family, when its series are measured anyway, and step 6 the later decisions on variantsand process changes}\DIFdelend \DIFaddbegin \DIFadd{, whose initial process study covers the released setting, supplies the floor but not the variants}\DIFaddend .

\subsection{Simulation and machine learning in this regime}

Within a design family the production record carried the decision\DIFaddbegin \DIFadd{, even against a rescaled simulation}\DIFaddend . The simulated force trend \DIFaddbegin \DIFadd{of the two draw-ins }\DIFaddend was 2.7 to 5.1 times too large, and calibrating the friction coefficient did not repair it, which points to the friction law and the blank-holder model rather than to a parameter. \DIFdelbegin \DIFdel{Learned models added an interval that is safe where the calibration alternatives are exchangeable with the new design and over-confident where they are not. }\DIFdelend A multi-task model cannot help a family with no produced alternative, because the correlation between families cannot be estimated without its data; \DIFdelbegin \DIFdel{the pooled Gaussian process, which carries the geometry as an input , is the borrowing of strength the data allow, and it }\DIFdelend \DIFaddbegin \DIFadd{in the pooled scopes the Gaussian process with a 0/1 geometry input is a restricted intrinsic coregionalisation model (equal variances, one positive correlation between the families, shared length scales), which }\DIFaddend improved coverage but not decisiveness. These findings concern one simulator of low fidelity\DIFdelbegin \DIFdel{for this production}\DIFdelend , used through offsets and one scale, and learners with six to nine produced \DIFdelbegin \DIFdel{settings each; a surrogate of these simulations inherits the simulator's distance, whatever its accuracy against the simulation. For the evaluation of }\DIFdelend \DIFaddbegin \DIFadd{alternatives at two or three force levels. For evaluating }\DIFaddend AI-based DFM support the lesson is methodological: \DIFdelbegin \DIFdel{held-out accuracy is not enough; }\DIFdelend a learned model should be compared with a production-record baseline, scored by evidence relation in a unit that production defines, and shown to its user with a guard band. \DIFdelbegin \DIFdel{Whether that makes designers' reliance match a rule's reliability \mbox{%DIFAUXCMD
\citep{lee2004trust}}\hskip0pt%DIFAUXCMD
, and when an abstention should go to an expert rather than to a trial \mbox{%DIFAUXCMD
\citep{madras2018predict}}\hskip0pt%DIFAUXCMD
, are hypotheses for a study with design }\DIFdelend \DIFaddbegin \DIFadd{A study with design teams could test two hypotheses on reliance \mbox{%DIFAUXCMD
\citep{lee2004trust}}\hskip0pt%DIFAUXCMD
, with cases of known production outcome: designers shown the guard-banded verdict accept a smaller share of wrong verdicts within the safe distance than designers shown the verdict alone, at an equal rate of right acceptances; and, as the same AI assistance helped some teams and hindered others \mbox{%DIFAUXCMD
\citep{zhang2021cautionary}}\hskip0pt%DIFAUXCMD
, the guard band narrows the spread of that share between }\DIFaddend teams.

\subsection{Implications for design research}

The \DIFdelbegin \DIFdel{framework gives several }\DIFdelend \DIFaddbegin \DIFadd{results give }\DIFaddend constructs of design research a production-referenced measure. For variant \DIFdelbegin \DIFdel{, adaptive }\DIFdelend and platform design \citep{pahl2007engineering, simpson2006product} \DIFdelbegin \DIFdel{, }\DIFdelend \DIFaddbegin \DIFadd{and the capability databases that serve them \mbox{%DIFAUXCMD
\citep{tata1999process, thornton2004variation}}\hskip0pt%DIFAUXCMD
, the results argue for indexing production records by their relation to }\DIFaddend the \DIFdelbegin \DIFdel{distances by evidence relation quantify how far production knowledge carries from produced variants to a new one, with the nearest produced setting as the baseline that any capability database \mbox{%DIFAUXCMD
\citep{thornton2004variation} }\hskip0pt%DIFAUXCMD
or learned model has to beat}\DIFdelend \DIFaddbegin \DIFadd{coming design rather than for a more sophisticated model on top}\DIFaddend . The safe distance is the part of a margin that covers predictive uncertainty, as distinct from the part that absorbs change \citep{eckert2017safety, eckert2019design}. \DIFdelbegin \DIFdel{For testing strategies \mbox{%DIFAUXCMD
\citep{thomke2001sequential, salado2018mathematical, camburn2017design}}\hskip0pt%DIFAUXCMD
, overlapped development \mbox{%DIFAUXCMD
\citep{krishnan1997model, terwiesch2002exchanging} }\hskip0pt%DIFAUXCMD
and }\DIFdelend \DIFaddbegin \DIFadd{The growth of the decisive distance from about 3 to 8 and over 30 floors says for which overlaps of development \mbox{%DIFAUXCMD
\citep{krishnan1997model, terwiesch2002exchanging} }\hskip0pt%DIFAUXCMD
preliminary manufacturability information can be released, the new-family safe distance, some 0.8\,mm of draw-in, how far a }\DIFaddend set-based design \citep{sobek1999toyota} \DIFdelbegin \DIFdel{, the distances say }\DIFdelend \DIFaddbegin \DIFadd{can narrow a feasible set before production evidence exists, and the distances and trial rates }\DIFaddend when a test is cheaper than a decision \DIFdelbegin \DIFdel{, when design-stage information is fit to act on and how wide a feasibility limit has to be. For validation, the evidence relation is the question of the application domain \mbox{%DIFAUXCMD
\citep{oberkampf2010verification}}\hskip0pt%DIFAUXCMD
, and the distances could serve as evidence of a risk-informed credibility assessment \mbox{%DIFAUXCMD
\citep{asme2018assessing}}\hskip0pt%DIFAUXCMD
}\DIFdelend \DIFaddbegin \DIFadd{\mbox{%DIFAUXCMD
\citep{thomke2001sequential, salado2018mathematical}}\hskip0pt%DIFAUXCMD
}\DIFaddend . In the terms of the design research methodology \citep{blessing2009drm}, the evaluation \DIFdelbegin \DIFdel{is a support evaluation: it }\DIFdelend shows structural validity\DIFdelbegin \DIFdel{in the sense of the validation square}\DIFdelend , that the metrics discriminate between rules and relations on real data \citep{seepersad2006validation}\DIFdelbegin \DIFdel{, and is a first step towards performance validity , which would require showing that decisions taken }\DIFdelend \DIFaddbegin \DIFadd{; performance validity would require better decisions }\DIFaddend with the framework\DIFdelbegin \DIFdel{are better than decisions taken without it}\DIFdelend .

\DIFaddbegin \begin{table}[t]
\tabsize
\setlength{\tabcolsep}{3pt}
\caption{What the framework requires, mapped onto other processes (an untested proposal)}\label{tab:processes}
\begin{tabular}{@{}L{1.9cm}L{2.4cm}L{2.5cm}L{2.7cm}L{2.7cm}@{}}
\toprule
Process & Batch and run & Main variation & Evidence relations & Floor \\
\midrule
Sheet forming (this paper) & consecutive parts of one series & drift in a run; coil, die state & variant, setting, family & batch centres in a run \\
Injection moulding & shots of one cavity and set-up & cavity, shot, set-up, material lot & cavity, process window, new mould & per cavity; cavity offsets as effects \\
Machining & consecutive parts between offsets & tool wear, offset-compensated & tool path, cutting data, new part & batch centres within an offset interval \\
Casting & castings of one pattern and melt & melt, pattern position & alloy, gating, new pattern & within a melt; melts as runs \\
Additive manufacturing & replicates in one build & build, position, orientation & orientation, settings, new part & within a build, via replicated anchors \\
\botrule
\end{tabular}
\end{table}

\DIFaddend \subsection{Reuse for other processes and AI-based DFM tools}\label{sec:reuse}

The framework needs several produced variants of a design family whose relations \DIFdelbegin \DIFdel{to each other }\DIFdelend span those of the coming designs, a series of consecutive parts or replicated batches of one design per variant, a design-stage predictor for every variant, and characteristics measured identically on parts and predictions\DIFdelbegin \DIFdel{. Here the long series served the floor, whereas calibration runs of }\DIFdelend \DIFaddbegin \DIFadd{; here rules calibrated on }\DIFaddend about 50 parts per variant \DIFdelbegin \DIFdel{qualified the rules }\DIFdelend \DIFaddbegin \DIFadd{scored }\DIFaddend as well as \DIFdelbegin \DIFdel{full series; how many long runsa floor needs was not tested. Table~\ref{tab:processes} maps these requirements onto the processes named most often in DFM; attribute }\DIFdelend \DIFaddbegin \DIFadd{with full series, and only the floor needed long runs. Attribute }\DIFaddend outcomes such as splits \DIFdelbegin \DIFdel{, porosity or short shots }\DIFdelend need a requirement on a rate\DIFdelbegin \DIFdel{rather than a quantile}\DIFdelend , which the framework does not yet cover. \DIFdelbegin \DIFdel{A predictor need not return an interval: any tool that returns meets, fails or trial for a stated requirement , such as a classifier at its operating threshold or a }\DIFdelend \DIFaddbegin \DIFadd{Other predictors need three adaptations. A classifier with a fixed threshold does not take the requirement as input, so its operating characteristic is estimated by binning test cases by their true margin, not by the exact estimator of Section~\ref{sec:distances}, unless the requirement is made an input; a probabilistic classifier needs two thresholds to return a trial. A }\DIFaddend language-model check \DIFdelbegin \DIFdel{given the requirement value, can be scored by sweeping the requirement , }\DIFdelend \DIFaddbegin \DIFadd{may return verdicts that are not monotone in the requirement and that vary between queries, so $\kappa$ }\DIFaddend and \DIFdelbegin \DIFdel{for }\DIFdelend \DIFaddbegin \DIFadd{$\omega$ are computed on a requirement grid with repeated queries, and the share of non-monotone verdict sequences is reported. For }\DIFaddend continuous design spaces the relation becomes a distance from the produced designs in descriptor space or the position relative to their hull\DIFaddbegin \DIFadd{, and each stratum needs about as many cases as a relation here: with 126 cases, some 240 verdicts, the 95\,\% threshold already rests on about a dozen. For a generative tool, its manufacturability check is scored}\DIFaddend . A minimal report comprises
\DIFaddbegin \begin{enumerate}
\item \DIFaddend the floor with its protocol\DIFdelbegin \DIFdel{and variance components , the }\DIFdelend \DIFaddbegin \DIFadd{, variance components and lag dependence;
}\item \DIFaddend decisive and safe distances by relation\DIFdelbegin \DIFdel{with refitting }\DIFdelend \DIFaddbegin \DIFadd{, with refit }\DIFaddend intervals and one-sided values\DIFdelbegin \DIFdel{, the coverage }\DIFdelend \DIFaddbegin \DIFadd{;
}\item \DIFadd{coverage and half-width }\DIFaddend of interval rules, \DIFaddbegin \DIFadd{and the source of their noise variance;
}\item \DIFaddend the nearest-produced-setting and no-simulation baselines\DIFdelbegin \DIFdel{, and }\DIFdelend \DIFaddbegin \DIFadd{;
}\item \DIFaddend the guard bands \DIFdelbegin \DIFdel{and guard performance . Distances per characteristic rest on as many cases as there are produced alternatives, so a team would pool characteristics with similar ratios of floor to part scatter, or accumulate evaluations over families}\DIFdelend \DIFaddbegin \DIFadd{with their prior, the trial rates and the performance of the guards;
}\item \DIFadd{for language-model tools, the prompt and the query protocol.
}\end{enumerate}
\DIFadd{Table~\ref{tab:processes}, an untested proposal, maps these requirements onto other processes; distances in floors are comparable within a protocol, not across processes, since a floor between builds, cavities or offset intervals is a between-unit quantity (here the floor between series is 1.5 to 9.1 times that within a run)}\DIFaddend . An attempt \DIFdelbegin \DIFdel{to apply the framework to an open dataset of }\DIFdelend \DIFaddbegin \DIFadd{on open }\DIFaddend powder bed fusion \DIFaddbegin \DIFadd{data }\DIFaddend of PA12 \citep{leirmo2021data} failed on the second requirement: \DIFdelbegin \DIFdel{grouping specimens by orientation class }\DIFdelend \DIFaddbegin \DIFadd{orientation classes }\DIFaddend mixed orientations in a batch, \DIFdelbegin \DIFdel{and eight identical anchor specimens }\DIFdelend \DIFaddbegin \DIFadd{eight identical anchors }\DIFaddend per build were too few for a floor \DIFdelbegin \DIFdel{; a between-build floor from replicated anchors would be the remedy}\DIFdelend \DIFaddbegin \DIFadd{within a build, and three builds give three differences between builds}\DIFaddend .

\DIFdelbegin %DIFDELCMD < \begin{table}[t]
%DIFDELCMD < \tabsize
%DIFDELCMD < \setlength{\tabcolsep}{3pt}
%DIFDELCMD < \caption{What the framework requires, mapped onto other processes}\label{tab:processes}
%DIFDELCMD < \begin{tabular}{@{}L{2.0cm}L{2.5cm}L{2.7cm}L{2.6cm}L{2.4cm}@{}}
%DIFDELCMD < \toprule
%DIFDELCMD < Process & Batch and run & Dominant variation & Evidence relations & Floor \\
%DIFDELCMD < \midrule
%DIFDELCMD < Sheet forming (this paper) & consecutive parts of one series & drift within a run; coil, die state & variant, setting, family & batch centres within a run \\
%DIFDELCMD < Injection moulding & shots of one cavity and set-up & cavity, shot, set-up & cavity, process window, new mould & within cavity, between cavities \\
%DIFDELCMD < Machining & consecutive parts between offsets & tool wear, compensated by offsets & tool path, cutting data, new part & between offset intervals \\
%DIFDELCMD < Casting & castings of one pattern and melt & melt, pattern position & alloy, gating, new pattern & between melts \\
%DIFDELCMD < Additive manufacturing & replicated specimens of one build & build, position, orientation & orientation, parameters, new part & between builds, via anchors \\
%DIFDELCMD < \botrule
%DIFDELCMD < \end{tabular}
%DIFDELCMD < \end{table}
%DIFDELCMD < 

%DIFDELCMD < %%%
\DIFdelend \subsection{Threats to validity}

The evaluation covers one material, one press and two geometries, \DIFdelbegin \DIFdel{and the }\DIFdelend \DIFaddbegin \DIFadd{with }\DIFaddend relations within a family \DIFdelbegin \DIFdel{come }\DIFdelend from a $3\times3$ factorial in which lubrication is mostly not resolved; the procedure is general, the numbers are not. The new family rests on two transfers between geometries that differ in three features at once, and the intervals of its distances hold the two fits fixed. Each alternative was produced once, so the floor covers drift within a run but not between days, coils or tool states, the reference $q_{95}$ is that of one run, and the order and dates of the series are unknown; between-series variation bounds the missing component at 1.2--9.8 floors\DIFdelbegin \DIFdel{(Section~\ref{sec:res-floor}). Scanner }\DIFdelend \DIFaddbegin \DIFadd{, and scanner }\DIFaddend drift cannot be separated from process drift. The extrapolation \DIFdelbegin \DIFdel{in force }\DIFdelend \DIFaddbegin \DIFadd{to 100\,kN }\DIFaddend is confounded with stroke speed, which the simulation \DIFdelbegin \DIFdel{does not represent, so the extrapolation results describe a force and speed change together}\DIFdelend \DIFaddbegin \DIFadd{omits}\DIFaddend . The requirements are one-sided limits on the 95th percentile of a characteristic averaged over sides, below the conformance of series release. Several rules, the grouped schemes\DIFdelbegin \DIFdel{and }\DIFdelend \DIFaddbegin \DIFadd{, }\DIFaddend the guard bands \DIFaddbegin \DIFadd{and their prior }\DIFaddend were added or chosen after earlier results had been seen and are evaluated on the same data, so the comparison between rules carries selection optimism \citep{cawley2010over}; \DIFdelbegin \DIFdel{the analysis was carried out by }\DIFdelend a single analyst \DIFaddbegin \DIFadd{carried out the analysis}\DIFaddend , with code and data open for replication.

%% ----------------------------------------------------------------------
\section{Conclusions}\label{sec:conclusion}

Design-stage manufacturability decisions can be evaluated against production on a scale \DIFdelbegin \DIFdel{that }\DIFdelend production defines: the decisive and safe distances in floors tell how close to a requirement a rule decides and how close it can be relied on not to mislead, and a guard band \DIFaddbegin \DIFadd{derived for the requirements a team expects }\DIFaddend turns them into a rule for \DIFdelbegin \DIFdel{a single decision. }%DIFDELCMD < 

%DIFDELCMD < %%%
\DIFdelend \DIFaddbegin \DIFadd{single decisions. }\DIFaddend On 18 deep-drawn \DIFdelbegin \DIFdel{design--process }\DIFdelend alternatives with matched simulations, the best attainable distance grew with \DIFdelbegin \DIFdel{the noveltyof the design}\DIFdelend \DIFaddbegin \DIFadd{novelty}\DIFaddend : 3.4 floors with a produced sibling \DIFdelbegin \DIFdel{, 8.5 }\DIFdelend \DIFaddbegin \DIFadd{(the nearest setting 0.85 for near-replicates, 4.9 for resolvable ones), 8.2 }\DIFaddend for a new process setting \DIFaddbegin \DIFadd{(4.7 interpolated; 4.9 extrapolated to 500\,kN and 8.9 to 100\,kN, with a stroke-speed change)}\DIFaddend , and no rule safe closer than \DIFdelbegin \DIFdel{16 }\DIFdelend \DIFaddbegin \DIFadd{16--17 }\DIFaddend floors for a new family\DIFaddbegin \DIFadd{, in either family's floor}\DIFaddend . Within a family the choice of rule mattered more than the relation\DIFdelbegin \DIFdel{, }\DIFdelend \DIFaddbegin \DIFadd{; }\DIFaddend release-level requirements lay inside every decisive distance \DIFdelbegin \DIFdel{, }\DIFdelend \DIFaddbegin \DIFadd{except a near-replicate's and need a trial without a produced sibling; }\DIFaddend and in this small-data regime the simulation as used, \DIFdelbegin \DIFdel{a model with a large and partly sign-inconsistent bias}\DIFdelend \DIFaddbegin \DIFadd{offset or rescaled}\DIFaddend , contributed only across families, while learned \DIFdelbegin \DIFdel{models added abstention rather than accuracy}\DIFdelend \DIFaddbegin \DIFadd{90\,\% intervals covered 85--96\,\% of truths with a produced sibling, under 50\,\% extrapolating or for a new family}\DIFaddend .

Four extensions follow: series \DIFdelbegin \DIFdel{produced on }\DIFdelend \DIFaddbegin \DIFadd{from }\DIFaddend several days, coils and tool states\DIFdelbegin \DIFdel{, so that the floor covers variation between runs; design }\DIFdelend \DIFaddbegin \DIFadd{; }\DIFaddend families with continuous geometric descriptors; a confirmatory \DIFdelbegin \DIFdel{evaluation }\DIFdelend \DIFaddbegin \DIFadd{test }\DIFaddend of the frozen procedure on \DIFdelbegin \DIFdel{another dataset or process}\DIFdelend \DIFaddbegin \DIFadd{other data}\DIFaddend ; and a study \DIFdelbegin \DIFdel{with design teams of whether decisions taken with the framework are better . The released code, which reproduces every result from the public data, is the starting point}\DIFdelend \DIFaddbegin \DIFadd{of whether design teams decide better with it}\DIFaddend .

%% ----------------------------------------------------------------------
\DIFaddbegin 

\DIFaddend \FloatBarrier
\begin{appendices}
\DIFaddbegin \renewcommand{\floatpagefraction}{0.4}
\DIFaddend \renewcommand{\theHtable}{\thetable}% distinct hyperlink targets for Tables A1 and A2

\DIFdelbegin %DIFDELCMD < \section{Notation, robustness and computation details}%%%
\DIFdelend \DIFaddbegin \section{Notation, further results and computation details}\DIFaddend \label{app:appendix}

Table~\ref{tab:notation} summarises the notation of Section~\ref{sec:method}, \DIFdelbegin \DIFdel{and }\DIFdelend Table~\ref{tab:robust} gives the distances under the robustness variants\DIFdelbegin \DIFdel{; the other results quoted in the text, such as guard bands, coverage}\DIFdelend , \DIFdelbegin \DIFdel{paired differences and distances per characteristic, are written by the released code and archived with it }\DIFdelend \DIFaddbegin \DIFadd{Table~\ref{tab:intervals} the coverage, width and centre of the interval rules, and Table~\ref{tab:step6} the procedure's decision rule by relation. Table~\ref{tab:mapping} traces every other number quoted in the text to the script that computes it and the result file that holds it in the archived release \mbox{%DIFAUXCMD
\citep{canseven2026archive}}\hskip0pt%DIFAUXCMD
}\DIFaddend .








\DIFdelbegin %DIFDELCMD < \begin{table}[h]
%DIFDELCMD < \tabsize
%DIFDELCMD < \setlength{\tabcolsep}{4pt}
%DIFDELCMD < \caption{Notation and terms}\label{tab:notation}
%DIFDELCMD < \begin{tabular}{@{}L{2.55cm}L{10.2cm}@{}}
%DIFDELCMD < \toprule
%DIFDELCMD < Symbol, term & Meaning \\
%DIFDELCMD < \midrule
%DIFDELCMD < $a$, $c$, $q_p(a,c)$ & alternative (design--process combination), characteristic, $p$-quantile of the parts of $a$ ($p=0.95$) \\
%DIFDELCMD < $R$, $d$ & requirement (upper limit on $c$ or $|c|$); its distance from the truth in floors, $R=q_p+dF_c$ \\
%DIFDELCMD < $[L_r,U_r]$, $V_r$ & interval of rule $r$ for $q_p$; verdict meets, fails or trial (Eq.~\ref{eq:verdict}) \\
%DIFDELCMD < $B$, $m$, $F_c$ & batch size (50); batch centre (10\,\% trimmed mean); production floor (Eq.~\ref{eq:floor}) \\
%DIFDELCMD < $\sigma_\mathrm{b}$, $\sigma_\mathrm{w}$ & standard deviation between batch means and between parts within a batch \\
%DIFDELCMD < ESR, MRC & effect-to-scatter ratio (Eq.~\ref{eq:esr}); minimum resolvable change, $F_c/|s_c|$ \\
%DIFDELCMD < $u$, $v$ & distances in floors by which an interval excludes the truth from above and below \\
%DIFDELCMD < $\kappa(\delta)$, $\omega(\delta)$ & shares of correct and wrong verdicts among the feasible verdicts at $|d|=\delta$ \\
%DIFDELCMD < $\delta_\mathrm{d}$, $\delta_\mathrm{s}$ & decisive distance ($\kappa\ge0.95$ beyond it) and safe distance ($\omega\le0.05$ beyond it) \\
%DIFDELCMD < FA, FR side & one-sided distances for failing alternatives (false accepts) and adequate ones (false rejects) \\
%DIFDELCMD < $g$ & guard band: observable margin beyond which at least 95\,\% of the decided verdicts were correct \\
%DIFDELCMD < Calibration alternatives & produced alternatives whose record calibrates a rule \\
%DIFDELCMD < Evidence relation & sibling (same setting produced), new process setting (interpolated, extrapolated), new family \\
%DIFDELCMD < Calibration budget & distances against the number $k$ of produced alternatives, by relation \\
%DIFDELCMD < Guard & family, range and novelty checks that send a design outside the calibration to trial \\
%DIFDELCMD < \botrule
%DIFDELCMD < \end{tabular}
%DIFDELCMD < \end{table}
%DIFDELCMD < 

%DIFDELCMD < %%%
%DIF < % BEGIN GENERATED robust
%DIFDELCMD < \begin{table}[h]
%DIFDELCMD < \tabsize
%DIFDELCMD < \setlength{\tabcolsep}{2.4pt}
%DIFDELCMD < \caption{Decisive/safe distances (floors) under the robustness variants}\label{tab:robust}
%DIFDELCMD < \begin{tabular}{@{}lrrrrrrrrr@{}}
%DIFDELCMD < \toprule
%DIFDELCMD <  & \multicolumn{4}{c}{New variant} & \multicolumn{3}{c}{New setting} & \multicolumn{2}{c}{New family} \\
%DIFDELCMD < \cmidrule(lr){2-5}\cmidrule(lr){6-8}\cmidrule(lr){9-10}
%DIFDELCMD < Variant & M1s & M2 & M5 & NN & M2 & M5 & NN & M2 & M5 \\
%DIFDELCMD < \midrule
%DIFDELCMD < Reference & 4.9 & 23/0.15 & 7.2/0.75 & 3.4 & 18/8.1 & 19/6.8 & 8.5 & 51/16 & 38/28 \\
%DIFDELCMD < Temperature-adjusted & 5.0 & 23/0.10 & 7.2/0.65 & 3.4 & 17/7.6 & 19/6.8 & 9.4 & 48/15 & 37/27 \\
%DIFDELCMD < First 150 parts dropped & 6.5 & 27/0.20 & 9.0/0.45 & 4.1 & 21/8.7 & 24/8.7 & 12 & 63/18 & 39/29 \\
%DIFDELCMD < Floor between series & 2.9 & 11/0 & 1.3/0.20 & 0.75 & 9.9/2.6 & 11/2.3 & 2.7 & 16/3.8 & 12/9.3 \\
%DIFDELCMD < Floor: batches of 25 & 4.6 & 21/0.15 & 6.4/0.70 & 2.9 & 16/6.9 & 17/6.3 & 7.7 & 43/14 & 31/23 \\
%DIFDELCMD < Floor: batches of 100 & 5.3 & 24/0.15 & 8.2/0.80 & 3.5 & 19/8.5 & 21/7.9 & 8.9 & 57/17 & 46/29 \\
%DIFDELCMD < Floor: 90\,\% quantile & 6.2 & 31/0.20 & 8.8/0.90 & 4.0 & 21/10 & 23/8.6 & 10 & 61/20 & 45/34 \\
%DIFDELCMD < Floor: 99\,\% quantile & 3.4 & 14/0.10 & 5.2/0.50 & 2.4 & 11/4.4 & 13/4.4 & 5.7 & 37/11 & 29/15 \\
%DIFDELCMD < Calibration series of 50 & 5.0 & 22/0.20 & 6.8/0.70 & 3.4 & 18/7.5 & 19/6.2 & 8.4 & 50/16 & 37/28 \\
%DIFDELCMD < Calibration series of 100 & 4.9 & 22/0.40 & 6.8/0.40 & 3.3 & 18/7.5 & 19/6.2 & 8.7 & 50/16 & 38/28 \\
%DIFDELCMD < Calibration series of 250 & 4.9 & 22/0.45 & 7.0/0.45 & 3.3 & 18/7.5 & 19/6.2 & 8.7 & 50/16 & 38/28 \\
%DIFDELCMD < Conformance $p=0.90$ & 4.9 & 23/0.20 & 7.2/0.70 & 3.3 & 18/7.5 & 18/6.4 & 8.4 & 51/15 & 38/28 \\
%DIFDELCMD < Conformance $p=0.99$ & 5.1 & 23/0 & 7.5/0.70 & 3.3 & 18/8.3 & 20/6.9 & 8.7 & 51/16 & 38/28 \\
%DIFDELCMD < Interval level 0.80 &  & 23/0.15 & 6.4/1.0 &  & 18/8.1 & 17/8.1 &  & 50/17 & 38/28 \\
%DIFDELCMD < Interval level 0.95 &  & 23/0.15 & 8.0/0.25 &  & 18/8.1 & 20/6.2 &  & 51/16 & 39/27 \\
%DIFDELCMD < Response surface & 4.8 & 25/0.30 & 7.2/0.70 & 3.4 & 18/7.6 & 20/6.9 & 8.5 & 49/25 & 37/27 \\
%DIFDELCMD < GP: sampling noise floor &  &  & 6.7/1.4 &  &  & 18/7.4 &  &  & 38/30 \\
%DIFDELCMD < GP: no noise floor &  &  & 6.7/1.5 &  &  & 18/7.4 &  &  & 38/30 \\
%DIFDELCMD < Unit: $q_{95}$ reproducibility & 5.3 & 17/0.15 & 6.9/0.40 & 2.8 & 15/6.0 & 18/6.2 & 6.9 & 41/13 & 34/30 \\
%DIFDELCMD < \midrule Refit bootstrap, $\delta_\mathrm{d}$ & 4.9--7.3 & 22--23 & 4.0--11 & 2.3--4.6 & 16--18 & 18--20 & 6.9--9.2 & 49--53 & 34--39 \\
%DIFDELCMD < \botrule
%DIFDELCMD < \end{tabular}
%DIFDELCMD < %%%
\footnotetext{\DIFdelFL{Rules refitted for every variant; empty cells: the variant does not affect the rule. Floor variants and the unit of $q_{95}$ reproducibility (95\,\% quantile of the difference between the $q_{95}$ of two batches of 100 parts) rescale the distances. Calibration series of $n$: the calibration alternatives' $q_{95}$ from their first $n$ parts. Conformance $p$: the decided statistic is the $p$-quantile. Response surface: nominal simulation and envelope from a quadratic fit over thickness and friction. Other kernels and descriptors of the GPs move M5 by at most 0.4 floors. Refit bootstrap: 95\,\% interval when the lubrication patterns of each geometry are resampled and every rule is refitted (200 resamples).}}
%DIFAUXCMD
%DIFDELCMD < \end{table}
%DIFDELCMD < %%%
%DIF < % END GENERATED robust
%DIFDELCMD < 

%DIFDELCMD < %%%
\DIFdelend \emph{Computation details.} Batches need at least 80\,\% valid values. M5 and M5n are Gaussian-process regressions (constant times anisotropic squared-exponential kernel plus white noise, normalised targets, \DIFaddbegin \DIFadd{hyperparameters by maximum marginal likelihood without hyperpriors, }\DIFaddend three restarts, length scales in $[0.1, 100]$) on a geometry indicator, the force in 100\,kN and the lubrication rank\DIFaddbegin \DIFadd{; the lower bound of the noise variance, the larger of the mean sampling variance of the calibration alternatives' $q_{95}$ and the pooled variance between their series of one geometry and force, is expressed on the normalised scale}\DIFaddend . NN takes the mean $q_{95}$ at the nearest force of the family, both neighbours when equally near\DIFaddbegin \DIFadd{; M1sn fits $q_{95}=a+b\,$force by least squares}\DIFaddend . Mc calibrates, per lubrication pattern, the friction coefficient whose simulated draw-ins are closest to the calibration centres. M3 and M3n use histogram gradient boosting (depth 3, 200 iterations, rate 0.08) on geometry, force, lubrication, sheet thickness and oil film, with 500 incoming conditions drawn from calibration parts of the same pattern; their jackknife+ interval \DIFdelbegin \DIFdel{combines }\DIFdelend \DIFaddbegin \DIFadd{takes }\DIFaddend the leave-one-out predictions for the new alternative\DIFdelbegin \DIFdel{with the absolute }\DIFdelend \DIFaddbegin \DIFadd{, shifted by the median }\DIFaddend leave-one-out \DIFdelbegin \DIFdel{residuals about their median}\DIFdelend \DIFaddbegin \DIFadd{residual, minus and plus the absolute residuals about that median, with its bounds at the extreme order statistics when $n$ is too small for the 0.9 quantiles}\DIFaddend . Guard bands lie on a grid of 0.25 floors. The refitting bootstrap draws the three lubrication patterns of each geometry with replacement 200 times, at least two of them distinct\DIFaddbegin \DIFadd{, and does not restrict the noise floor to distinct series}\DIFaddend . Software: Python 3.11 with NumPy, pandas, SciPy, scikit-learn \citep{pedregosa2011scikit} and Matplotlib.
\DIFaddbegin 

\begin{table}[!ht]
\tabsize
\setlength{\tabcolsep}{4pt}
\caption{Notation and terms}\label{tab:notation}
\begin{tabular}{@{}L{2.55cm}L{10.2cm}@{}}
\toprule
Symbol, term & Meaning \\
\midrule
$a$, $c$, $q_p(a,c)$ & alternative (design--process combination), characteristic, $p$-quantile of the parts of $a$ ($p=0.95$) \\
$R$, $d$ & requirement (upper limit on $c$ or $|c|$); its distance from the truth in floors, $R=q_p+dF_c$ \\
$[L_r,U_r]$, $V_r$ & interval of rule $r$ for $q_p$; verdict meets, fails or trial (Eq.~\ref{eq:verdict}) \\
$B$, $m$, $F_c$ & batch size (50); batch centre (10\,\% trimmed mean); production floor (Eq.~\ref{eq:floor}) \\
$\sigma_\mathrm{b}$, $\sigma_\mathrm{w}$ & standard deviation between batch means and between parts within a batch \\
ESR, MRC & effect-to-scatter ratio (Eq.~\ref{eq:esr}); minimum resolvable change, $F_c/|s_c|$ \\
$u$, $v$ & distances in floors by which an interval excludes the truth from above and below \\
$\kappa(\delta)$, $\omega(\delta)$ & shares of correct and wrong verdicts among the feasible verdicts at $|d|=\delta$ \\
$\delta_\mathrm{d}$, $\delta_\mathrm{s}$ & decisive distance ($\kappa\ge0.95$ beyond it) and safe distance ($\omega\le0.05$ beyond it) \\
FA, FR side & one-sided distances for failing alternatives (false accepts) and adequate ones (false rejects) \\
$g$ & guard band: observable margin beyond which at least 95\,\% of the decided verdicts were correct, for a stated prior of requirement distances, here the 47 grid distances within $\pm20$ floors with equal weights (Section~\ref{sec:procedure}); not a property of rule and relation alone \\
Calibration alternatives & produced alternatives whose record calibrates a rule \\
Evidence relation & sibling (same setting produced), new process setting (interpolated, extrapolated), new family \\
Calibration budget & distances against the number $k$ of produced alternatives, by relation \\
Guard & family, range and novelty checks that send a design outside the calibration to trial \\
Measurement resolution & step between adjacent values that the measuring system reports (gauge sense) \\
\botrule
\end{tabular}
\end{table}

%DIF > % BEGIN GENERATED robust
\begin{table}[!ht]
\tabsize
\setlength{\tabcolsep}{1.6pt}
\caption{Decisive/safe distances (floors) under the robustness variants}\label{tab:robust}
\begin{tabular}{@{}lrrrrrrrrrr@{}}
\toprule
 & \multicolumn{4}{c}{New variant} & \multicolumn{4}{c}{New setting} & \multicolumn{2}{c}{New family} \\
\cmidrule(lr){2-5}\cmidrule(lr){6-9}\cmidrule(lr){10-11}
Variant & M1s & M2 & M5 & NN & M2 & M5 & NN & M1sn & M2 & M5 \\
\midrule
Reference & 4.9 & 23/0.15 & 7.2/0.75 & 3.4 & 18/8.1 & 19/6.8 & 8.5 & 8.2 & 51/16 & 38/28 \\
Temperature-adjusted & 5.0 & 23/0.10 & 7.2/0.65 & 3.4 & 17/7.6 & 19/6.8 & 9.4 & 8.9 & 48/15 & 37/27 \\
Warm-up dropped & 6.5 & 27/0.20 & 9.0/0.45 & 4.1 & 21/8.7 & 24/8.7 & 12 & 8.8 & 63/18 & 39/29 \\
Floor between series & 2.9 & 11/0 & 1.3/0.20 & 0.75 & 9.9/2.6 & 11/2.3 & 2.7 & 4.5 & 16/3.8 & 12/9.3 \\
Floor: batches of 25 & 4.6 & 21/0.15 & 6.4/0.70 & 2.9 & 16/6.9 & 17/6.3 & 7.7 & 7.3 & 43/14 & 31/23 \\
Floor: batches of 100 & 5.3 & 24/0.15 & 8.2/0.80 & 3.5 & 19/8.5 & 21/7.9 & 8.9 & 8.9 & 57/17 & 46/29 \\
Floor: 90\,\% quantile & 6.2 & 31/0.20 & 8.8/0.90 & 4.0 & 21/10 & 23/8.6 & 10 & 10 & 61/20 & 45/34 \\
Floor: 99\,\% quantile & 3.4 & 14/0.10 & 5.2/0.50 & 2.4 & 11/4.4 & 13/4.4 & 5.7 & 6.1 & 37/11 & 29/15 \\
Calibration: 50 parts & 5.0 & 22/0.20 & 6.8/0.70 & 3.4 & 18/7.5 & 19/6.2 & 8.4 & 7.7 & 50/16 & 37/28 \\
Calibration: 100 parts & 4.9 & 22/0.40 & 6.8/0.40 & 3.3 & 18/7.5 & 19/6.2 & 8.7 & 7.8 & 50/16 & 38/28 \\
Calibration: 250 parts & 4.9 & 22/0.45 & 7.0/0.45 & 3.3 & 18/7.5 & 19/6.2 & 8.7 & 8.3 & 50/16 & 38/28 \\
Conformance $p=0.90$ & 4.9 & 23/0.20 & 7.2/0.70 & 3.3 & 18/7.5 & 18/6.4 & 8.4 & 8.1 & 51/15 & 38/28 \\
Conformance $p=0.99$ & 5.1 & 23/0 & 7.5/0.70 & 3.3 & 18/8.3 & 20/6.9 & 8.7 & 8.4 & 51/16 & 38/28 \\
Interval level 0.80 &  & 23/0.15 & 6.4/1.0 &  & 18/8.1 & 17/8.1 &  &  & 50/17 & 38/28 \\
Interval level 0.95 &  & 23/0.15 & 8.0/0.25 &  & 18/8.1 & 20/6.2 &  &  & 51/16 & 39/27 \\
Response surface & 4.8 & 25/0.30 & 7.2/0.70 & 3.4 & 18/7.6 & 20/6.9 & 8.5 & 8.2 & 49/25 & 37/27 \\
GP: sampling bound &  &  & 6.7/1.4 &  &  & 18/7.4 &  &  &  & 38/30 \\
Unit: $q_{95}$ floor & 5.3 & 17/0.15 & 6.9/0.40 & 2.8 & 15/6.0 & 18/6.2 & 6.9 & 7.1 & 41/13 & 34/30 \\
\midrule Refit, $\delta_\mathrm{d}$ & 4.9--7.3 & 22--23 & 4.0--11 & 2.3--4.6 & 16--18 & 18--20 & 6.9--9.2 & 7.3--10 & 49--53 & 34--39 \\
\botrule
\end{tabular}
\footnotetext{\DIFaddFL{Rules refitted for every variant; empty cells: the variant does not affect the rule; a new family in its own floor. The floor between series contains the lubrication effect and, for a new variant, the series that NN averages, which makes it close to circular there. Warm-up: first 150 parts. Calibration: calibration $q_{95}$ from the first $n$ parts, truths and floors from the full series. Response surface: simulations from a quadratic fit over thickness and friction; other GP kernels and descriptors move M5 by at most 0.4 floors. Unit: the $q_{95}$ floor, 95\,\% quantile of the difference between the $q_{95}$ of two batches of 100 parts. Refit: 95\,\% interval of the decisive distance over 200 draws of each geometry's lubrication patterns with replacement, every rule refitted (Section~\ref{sec:evaldesign}).}}
\end{table}
%DIF > % END GENERATED robust

%DIF > % BEGIN GENERATED intervals
\begin{table}[!ht]
\tabsize
\setlength{\tabcolsep}{2.6pt}
\caption{Interval rules by relation: coverage, half-width and the interval centre as a point rule}\label{tab:intervals}
\begin{tabular}{@{}lrrrrrrrrrrrr@{}}
\toprule
 & \multicolumn{3}{c}{Sibling produced} & \multicolumn{3}{c}{Interpolated setting} & \multicolumn{3}{c}{Extrapolated setting} & \multicolumn{3}{c}{New family} \\
\cmidrule(lr){2-4}\cmidrule(lr){5-7}\cmidrule(lr){8-10}\cmidrule(lr){11-13}
Rule & Cov & HW & Ctr & Cov & HW & Ctr & Cov & HW & Ctr & Cov & HW & Ctr \\
\midrule
M2 & 89 & 7.6 & 8.6 & 71 & 5.7 & 10 & 44 & 4.8 & 13 & 47 & 7.6 & 34 \\
M2n & 90 & 5.9 & 9.1 & 93 & 5.0 & 4.8 & 33 & 3.1 & 12 & 12 & 5.7 & 131 \\
M3 & 96 & 9.3 & 9.9 & 43 & 7.0 & 19 & 43 & 8.9 & 21 & 40 & 9.7 & 29 \\
M3n & 89 & 3.0 & 3.7 & 36 & 3.4 & 11 & 48 & 2.3 & 9.9 & 7 & 3.3 & 130 \\
M5 & 86 & 2.2 & 3.3 & 71 & 6.4 & 9.8 & 46 & 4.6 & 13 & 12 & 2.3 & 34 \\
M5n & 85 & 2.3 & 3.1 & 93 & 5.8 & 4.7 & 36 & 3.7 & 12 & 0 & 2.2 & 130 \\
\botrule
\end{tabular}
\footnotetext{\DIFaddFL{Cov: share of cases whose true $q_{95}$ lies in the interval (\%; nominal 90). HW: median half-width (floors). Ctr: decisive distance of the interval centre used as a point rule (floors). A new family in the floor of the produced family. The GP noise variance sits at its lower bound in 97 and 99\,\% of the M5 and M5n fits with a sibling and in every new-setting and new-family fit.}}
\end{table}
%DIF > % END GENERATED intervals

%DIF > % BEGIN GENERATED step6
\begin{table}[!ht]
\tabsize
\setlength{\tabcolsep}{2.2pt}
\caption{The procedure's decision rule (step 6) by relation: guard band, distances and trial shares}\label{tab:step6}
\begin{tabular}{@{}lrrrrrrrrrrrr@{}}
\toprule
 & \multicolumn{4}{c}{Sibling produced} & \multicolumn{4}{c}{Interpolated setting} & \multicolumn{4}{c}{Extrapolated setting} \\
\cmidrule(lr){2-5}\cmidrule(lr){6-9}\cmidrule(lr){10-13}
Rule & $g$ & $\delta_\mathrm{d}/\delta_\mathrm{s}$ & $T_{10}$ & $T_{20}$ & $g$ & $\delta_\mathrm{d}/\delta_\mathrm{s}$ & $T_{10}$ & $T_{20}$ & $g$ & $\delta_\mathrm{d}/\delta_\mathrm{s}$ & $T_{10}$ & $T_{20}$ \\
\midrule
M1 & none & -- & 100 & 100 & 8.75 & 20/2.1 & 44 & 4 & none & -- & 100 & 100 \\
M1s & 2.0 & 6.9/2.9 & 1 & 0 & none & -- & 100 & 100 & none & -- & 100 & 100 \\
M2 & 0 & 23/0.15 & 33 & 8 & 0.75 & 19/2.0 & 31 & 4 & none & -- & 100 & 100 \\
M5 & 0 & 7.2/0.75 & 1 & 0 & 0.5 & 20/2.3 & 40 & 4 & none & -- & 100 & 100 \\
M5n & 0 & 7.2/0.55 & 0 & 0 & 0 & 16/0 & 25 & 0 & 6.75 & 25/1.9 & 48 & 15 \\
NN & 0.25 & 3.6/3.1 & 0 & 0 & 2.0 & 6.7/2.7 & 0 & 0 & 7.0 & 17/2.5 & 26 & 1 \\
M3 & 0 & 23/0 & 45 & 6 & 17.75 & 47/0 & 94 & 69 & 16.75 & 51/0.05 & 85 & 69 \\
M3n & 0 & 11/0.05 & 9 & 0 & 6.0 & 21/1.0 & 55 & 6 & 6.5 & 19/1.3 & 47 & 3 \\
M1sn & 0.75 & 4.3/2.8 & 0 & 0 & 2.0 & 6.7/2.7 & 0 & 0 & 6.0 & 15/2.9 & 21 & 0 \\
\botrule
\end{tabular}
\footnotetext{\DIFaddFL{$g$: guard band (floors) under the reference prior (Section~\ref{sec:procedure}); $\delta_\mathrm{d}/\delta_\mathrm{s}$: distances of the rule with its interval widened by $g$; $T_{10}$, $T_{20}$: share of verdicts sent to trial (\%) at requirements 10 and 20 floors from the truth. None: no band up to 20 floors reaches 95\,\%, so the procedure sends every design to trial; so also M0 and M0w in every relation, Mc with a sibling and for an extrapolated setting (for an interpolated one $g=18$), and every rule for a new family, scored in the produced family's floor.}}
\end{table}
%DIF > % END GENERATED step6

\begin{table}[!ht]
\footnotesize
\setlength{\tabcolsep}{3pt}
\caption{Where the numbers quoted in the text come from: scripts and result files (\texttt{results/*.csv})}\label{tab:mapping}
\begin{tabular}{@{}L{1.0cm}L{4.1cm}L{7.2cm}@{}}
\toprule
Section & Results & Script: result files \\
\midrule
3.4, 4.1, 4.3 & failing-side weight, run metadata and oil film, scan-line scatter, measurement resolution & panel.py: panel\_failing\_weight, panel\_runs, panel\_resolution\_steps; measurement.py: meas\_scan \\
5.1 & time order, lag profile, normal model, subgroups, part-to-part share, between series, $q_{95}$ reproducibility, in-line drift & alternatives.py: alt\_floor; sensitivity.py: sens\_drift; panel.py: panel\_variogram, panel\_floor\_protocol; measurement.py: meas\_floor; robustness.py: rob\_floor\_between; inline.py: inline\_drift \\
5.2 & resolvable contrasts and changes, simulation offsets, force effect, arm sign, cup depth & alternatives.py: alt\_effects, alt\_mrc; robustness.py: rob\_bhf\_pairs; decisions.py: dec\_alternatives; panel.py: panel\_force\_effect, panel\_arm\_sign, panel\_m0\_split \\
5.3--5.4 & distances in both floors, paired differences, strata, force levels, decomposition, pooled coverage & decisions.py: dec\_resolution, dec\_paired, dec\_coverage; panel.py: panel\_source\_floor(\_paired), panel\_sibling\_strata, panel\_force\_levels, panel\_decomposition \\
5.5 & simulation pairs, drawing nominal, surrogate, coverage, GP fits & decisions.py: dec\_paired, dec\_coverage, dec\_gp; panel.py: panel\_nominal\_offset, panel\_source\_floor(\_paired); design\_space.py: ds\_surrogate(\_rule) \\
5.6, 5.9 & budget, short series, robustness variants, refit, $q_{95}$ unit & budget.py: budget\_design, budget\_resolution; robustness.py: rob\_decisions, rob\_refit(\_paired); panel.py: panel\_q95\_unit \\
5.7 & margins, guard bands and priors, guards, worked example & panel.py: panel\_step6(\_curves), panel\_guard\_priors; guard.py: dec\_guard; decisions.py: dec\_intervals \\
5.8 & capability-anchored limits, ISO 2768-1 scenario, costs & panel.py: panel\_capability; scenarios.py: scen\_strata, scen\_scores, scen\_cost\_map \\
6.5 & PA12 attempt & case2\_pa12.py: case2\_floor, case2\_effects \\
\botrule
\end{tabular}
\footnotetext{\texttt{\DIFaddFL{scripts/run\_all.sh}} \DIFaddFL{runs the scripts in order; }\texttt{\DIFaddFL{make\_tables.py}} \DIFaddFL{writes every table and }\texttt{\DIFaddFL{figs.py}} \DIFaddFL{every figure from the result files.}}
\end{table}

\DIFaddend \end{appendices}

\backmatter

\bmhead{Acknowledgements}
The author thanks the authors of the RDDAC and DDACS datasets at the University of Stuttgart for publishing the data on which this work rests.

\section*{Declarations}

\bmhead{Funding}
No specific funding was received for this work.

\bmhead{Competing interests}
The author declares no competing interests.

\bmhead{Ethics approval and consent}
Not applicable.

\bmhead{Data availability}
All data are public. RDDAC: doi:10.18419/DARUS-5589 (version 2.0); DDACS: doi:10.18419/DARUS-4801 (version 3.0); both licensed CC BY 4.0. The extracted per-part and per-simulation feature tables \DIFdelbegin \DIFdel{are provided }\DIFdelend \DIFaddbegin \DIFadd{and all result files are archived }\DIFaddend with the code \DIFaddbegin \DIFadd{\mbox{%DIFAUXCMD
\citep{canseven2026archive}}\hskip0pt%DIFAUXCMD
}\DIFaddend .

\bmhead{Code availability}
The complete analysis pipeline\DIFdelbegin \DIFdel{is released under an open licence at }%DIFDELCMD < \url{https://github.com/htcanseven/Claude} %%%
\DIFdel{(directory }\texttt{\DIFdel{scripts/}}%DIFAUXCMD
\DIFdel{) and }\DIFdelend \DIFaddbegin \DIFadd{, which }\DIFaddend reproduces every number in this paper from the public data\DIFdelbegin \DIFdel{; a versioned archive of the code and }\DIFdelend \DIFaddbegin \DIFadd{, is archived with }\DIFaddend the feature tables \DIFdelbegin \DIFdel{is deposited on Zenodo \pending{DOI}}\DIFdelend \DIFaddbegin \DIFadd{and result files of this version on Zenodo \mbox{%DIFAUXCMD
\citep{canseven2026archive}}\hskip0pt%DIFAUXCMD
; Table~\ref{tab:mapping} maps the numbers to its scripts and files. A working copy is at }\url{https://github.com/htcanseven/Claude} \DIFadd{(directory }\texttt{\DIFadd{scripts/}}\DIFadd{)}\DIFaddend .

\bmhead{Author contributions}
H.T.C. conceived the method, designed and carried out the analysis, and wrote the paper.

\bmhead{Use of generative AI}
A large language model (Claude, Anthropic) was used under the author's direction for code development, literature retrieval, data analysis, drafting and editing, and simulated peer review of earlier versions. The author designed the study, verified all analyses, checked every reference against its published record, and takes full responsibility for the content. All figures are produced by the released code from the data; no AI-generated images are used.


\DIFdelbegin %DIFDELCMD < \FloatBarrier
%DIFDELCMD < %%%
\DIFdelend %% ----------------------------------------------------------------------
\DIFaddbegin \FloatBarrier
\DIFaddend {\small
\bibliography{refs}}

\end{document}
